\documentclass[11pt,letterpaper]{article}
\usepackage[utf8]{inputenc}
\usepackage[T1]{fontenc}
\usepackage[spanish,es-noshorthands,es-tabla]{babel}
\usepackage{lmodern}
\usepackage[margin=2.2cm]{geometry}
\usepackage{amsmath,amssymb,amsthm,mathtools}
\usepackage{enumitem}
\usepackage{booktabs,array,tabularx}
\usepackage{multicol}
\usepackage{xcolor}
\usepackage[most]{tcolorbox}
\usepackage{tikz}
\usetikzlibrary{shapes.geometric,arrows.meta,positioning,calc}
\usepackage{fancyhdr}
\usepackage{titlesec}
\usepackage{listings}
\usepackage[hidelinks]{hyperref}

\definecolor{azul}{HTML}{1F4E79}
\definecolor{azulclaro}{HTML}{E8F0F8}
\definecolor{verde}{HTML}{2E7D32}
\definecolor{verdeclaro}{HTML}{EAF5EA}
\definecolor{naranja}{HTML}{B85C00}
\definecolor{naranjaclaro}{HTML}{FFF3E0}
\definecolor{gris}{HTML}{555555}
\definecolor{rojo}{HTML}{B71C1C}

\titleformat{\section}{\Large\bfseries\color{azul}}{\thesection}{0.6em}{}
\titleformat{\subsection}{\large\bfseries\color{azul}}{\thesubsection}{0.5em}{}
\titleformat{\subsubsection}{\normalsize\bfseries\color{gris}}{\thesubsubsection}{0.5em}{}

\setlength{\parindent}{0pt}
\setlength{\parskip}{4pt}
\setlist{itemsep=2pt,topsep=3pt}

\pagestyle{fancy}
\fancyhf{}
\renewcommand{\headrulewidth}{0.4pt}

% ---------- Entornos ----------
\newcounter{ej}
\newtcolorbox{ejercicio}[2][]{enhanced,breakable,colback=white,colframe=azul,
  boxrule=0.6pt,arc=2pt,left=6pt,right=6pt,top=4pt,bottom=4pt,
  title={\textbf{#2}},fonttitle=\small\bfseries,coltitle=white,colbacktitle=azul,#1}
\newtcolorbox{solucion}[2][]{enhanced,breakable,colback=verdeclaro,colframe=verde,
  boxrule=0.6pt,arc=2pt,left=6pt,right=6pt,top=4pt,bottom=4pt,
  title={\textbf{#2}},fonttitle=\small\bfseries,coltitle=white,colbacktitle=verde,#1}
\newtcolorbox{nota}[1][]{enhanced,breakable,colback=naranjaclaro,colframe=naranja,
  boxrule=0.5pt,arc=2pt,left=6pt,right=6pt,top=3pt,bottom=3pt,#1}
\newtcolorbox{formulario}[1][]{enhanced,breakable,colback=azulclaro,colframe=azul,
  boxrule=0.5pt,arc=2pt,left=6pt,right=6pt,top=3pt,bottom=3pt,#1}

\newcommand{\dif}[1]{%
  \ifcase#1\or\textcolor[HTML]{66BB6A}{$\bigstar$}\or\textcolor[HTML]{FFA726}{$\bigstar\bigstar$}\or\textcolor[HTML]{EF5350}{$\bigstar\bigstar\bigstar$}\fi}
\newcommand{\fuente}[1]{\hfill{\footnotesize\textcolor{white!80!azul}{[#1]}}}
\newcommand{\ptst}[1]{{\footnotesize\textcolor{white!80!azul}{(#1 pts)}}}
\newcommand{\pts}[1]{\textcolor{gris}{(#1 pts)}}

% ---------- Probabilidad y RL ----------
\newcommand{\indep}{\perp\!\!\!\perp}
\newcommand{\nindep}{\not\!\perp\!\!\!\perp}
\newcommand{\argmax}{\operatorname*{arg\,max}}
\newcommand{\argmin}{\operatorname*{arg\,min}}
\newcommand{\E}{\mathbb{E}}
\newcommand{\R}{\mathbb{R}}
\tikzset{
  bn/.style={draw,circle,minimum size=0.8cm,inner sep=0pt,font=\small\bfseries,fill=azulclaro},
  bno/.style={bn,fill=naranjaclaro},
  barc/.style={-{Stealth[length=2.2mm]},thick},
  st/.style={draw,rounded corners,minimum width=1.0cm,minimum height=0.8cm,font=\small\bfseries,fill=azulclaro},
  term/.style={st,fill=verdeclaro},
}
% ---------- Lógica ----------
\newcommand{\imp}{\Rightarrow}
\newcommand{\sii}{\Leftrightarrow}
\newcommand{\ent}{\models}
\newcommand{\der}{\vdash}
\newcommand{\p}[1]{\ensuremath{\textit{#1}}}

% ---------- Árboles de juego ----------
\tikzset{
  maxn/.style={draw,regular polygon,regular polygon sides=3,minimum size=0.75cm,inner sep=0pt,fill=azulclaro},
  minn/.style={draw,regular polygon,regular polygon sides=3,shape border rotate=180,minimum size=0.75cm,inner sep=0pt,fill=naranjaclaro},
  chn/.style={draw,circle,minimum size=0.6cm,inner sep=0pt,fill=verdeclaro},
  hoja/.style={draw,rectangle,minimum size=0.5cm,inner sep=1pt,font=\small},
  gnode/.style={draw,circle,minimum size=0.75cm,inner sep=0pt,font=\small\bfseries},
  garc/.style={-{Stealth[length=2mm]},thick},
}

\lstset{basicstyle=\ttfamily\small,columns=fullflexible,frame=single,framerule=0.3pt,rulecolor=\color{gris},
  literate={á}{{\'a}}1 {é}{{\'e}}1 {í}{{\'i}}1 {ó}{{\'o}}1 {ú}{{\'u}}1 {ñ}{{\~n}}1}

% ================= Figuras compartidas (guía y respuestas) — Examen 2 =================

% ---- Red del estudiante ----
\newcommand{\figEst}{%
\begin{tikzpicture}[scale=0.95]
  \node[bn] (D) at (0,2) {D};
  \node[bn] (I) at (2.4,2) {I};
  \node[bn] (C) at (1.2,0.8) {C};
  \node[bn] (E) at (3.6,0.8) {E};
  \node[bn] (R) at (1.2,-0.4) {R};
  \draw[barc] (D) -- (C); \draw[barc] (I) -- (C); \draw[barc] (I) -- (E); \draw[barc] (C) -- (R);
\end{tikzpicture}}

\newcommand{\tabEst}{%
{\small
\begin{tabular}{c}
$P(+d)=0.4$ \qquad $P(+i)=0.3$\\[3pt]
\begin{tabular}{cc|c}
$D$ & $I$ & $P(+c\mid D,I)$\\\hline
$+d$ & $+i$ & 0.7\\ $+d$ & $-i$ & 0.1\\ $-d$ & $+i$ & 0.9\\ $-d$ & $-i$ & 0.4
\end{tabular}\quad
\begin{tabular}{c|c}
$I$ & $P(+e\mid I)$\\\hline $+i$ & 0.8\\ $-i$ & 0.1
\end{tabular}\quad
\begin{tabular}{c|c}
$C$ & $P(+r\mid C)$\\\hline $+c$ & 0.9\\ $-c$ & 0.2
\end{tabular}
\end{tabular}}}

% ---- Grafo de d-separación (8 nodos) ----
\newcommand{\figDsep}{%
\begin{tikzpicture}[scale=0.95]
  \node[bn] (A) at (0,3) {A};
  \node[bn] (B) at (2.4,3) {B};
  \node[bn] (C) at (1.2,1.8) {C};
  \node[bn] (D) at (3.6,1.8) {D};
  \node[bn] (E) at (1.2,0.6) {E};
  \node[bn] (F) at (3.6,0.6) {F};
  \node[bn] (G) at (2.4,-0.6) {G};
  \node[bn] (H) at (2.4,-1.8) {H};
  \draw[barc] (A)--(C); \draw[barc] (B)--(C); \draw[barc] (B)--(D); \draw[barc] (C)--(E);
  \draw[barc] (D)--(F); \draw[barc] (E)--(G); \draw[barc] (F)--(G); \draw[barc] (G)--(H);
\end{tikzpicture}}

% ---- Red de alarma ----
\newcommand{\figAlarma}{%
\begin{tikzpicture}[scale=0.9]
  \node[bn] (B) at (0,1.2) {B}; \node[bn] (E) at (2.4,1.2) {E};
  \node[bn] (A) at (1.2,0) {A};
  \node[bn] (J) at (0,-1.2) {J}; \node[bn] (M) at (2.4,-1.2) {M};
  \draw[barc] (B)--(A); \draw[barc] (E)--(A); \draw[barc] (A)--(J); \draw[barc] (A)--(M);
\end{tikzpicture}}

% ---- Red del aspersor ----
\newcommand{\figAsp}{%
\begin{tikzpicture}[scale=0.9]
  \node[bn] (C) at (1.2,1.2) {C};
  \node[bn] (S) at (0,0) {S}; \node[bn] (R) at (2.4,0) {R};
  \node[bn] (W) at (1.2,-1.2) {W};
  \draw[barc] (C)--(S); \draw[barc] (C)--(R); \draw[barc] (S)--(W); \draw[barc] (R)--(W);
\end{tikzpicture}}

\newcommand{\tabAsp}{%
{\small
\begin{tabular}{c}
$P(+c)=0.5$\qquad $P(+s\mid+c)=0.1$,\ $P(+s\mid-c)=0.5$\qquad $P(+r\mid+c)=0.8$,\ $P(+r\mid-c)=0.2$\\[3pt]
\begin{tabular}{cc|cccc}
$(S,R)$ & & $(+s,+r)$ & $(+s,-r)$ & $(-s,+r)$ & $(-s,-r)$\\\hline
$P(+w\mid S,R)$ & & 0.99 & 0.90 & 0.90 & 0.01
\end{tabular}
\end{tabular}}}

% ---- Red para el orden de eliminación ----
\newcommand{\figOrden}{%
\begin{tikzpicture}[scale=0.9]
  \node[bn] (A) at (0,1.2) {A}; \node[bn] (B) at (2.4,1.2) {B};
  \node[bn] (C) at (1.2,0) {C};
  \node[bn] (D) at (0,-1.2) {D}; \node[bn] (E) at (2.4,-1.2) {E};
  \node[bno] (F) at (1.2,-2.4) {F};
  \draw[barc] (A)--(C); \draw[barc] (B)--(C); \draw[barc] (C)--(D); \draw[barc] (C)--(E);
  \draw[barc] (D)--(F); \draw[barc] (E)--(F);
\end{tikzpicture}}

% ---- Red del servidor ----
\newcommand{\figSrv}{%
\begin{tikzpicture}[scale=0.9]
  \node[bn] (A) at (0,1.2) {A}; \node[bn] (U) at (2.4,1.2) {U};
  \node[bn] (C) at (1.2,0) {C};
  \node[bn] (M) at (0,-1.2) {M}; \node[bn] (Q) at (2.4,-1.2) {Q};
  \draw[barc] (A)--(C); \draw[barc] (U)--(C); \draw[barc] (C)--(M); \draw[barc] (C)--(Q);
\end{tikzpicture}}

\newcommand{\tabSrv}{%
{\small
\begin{tabular}{c}
$P(+a)=0.05$\qquad $P(+u)=0.1$\\[3pt]
\begin{tabular}{cc|c}
$A$ & $U$ & $P(+c\mid A,U)$\\\hline
$+a$ & $+u$ & 0.95\\ $+a$ & $-u$ & 0.80\\ $-a$ & $+u$ & 0.60\\ $-a$ & $-u$ & 0.05
\end{tabular}\quad
\begin{tabular}{c|cc}
$C$ & $P(+m\mid C)$ & $P(+q\mid C)$\\\hline $+c$ & 0.9 & 0.7\\ $-c$ & 0.1 & 0.05
\end{tabular}
\end{tabular}}}

% ---- Red del automóvil (simulacro 1) ----
\newcommand{\figCar}{%
\begin{tikzpicture}[scale=0.9]
  \node[bn] (B) at (0,1.2) {B}; \node[bn] (F) at (2.4,1.2) {F};
  \node[bn] (S) at (1.2,0) {S}; \node[bn] (L) at (-1.2,0) {L};
  \draw[barc] (B)--(S); \draw[barc] (F)--(S); \draw[barc] (B)--(L);
\end{tikzpicture}}

% ---- Red del simulacro 2 ----
\newcommand{\figDdos}{%
\begin{tikzpicture}[scale=0.9]
  \node[bn] (N) at (1.2,1.2) {N};
  \node[bn] (R) at (0,0) {R}; \node[bn] (S) at (2.4,0) {S};
  \node[bn] (W) at (1.2,-1.2) {W}; \node[bn] (T) at (-1.2,-1.2) {T};
  \draw[barc] (N)--(R); \draw[barc] (N)--(S); \draw[barc] (R)--(W); \draw[barc] (S)--(W); \draw[barc] (R)--(T);
\end{tikzpicture}}

% ---- Corredor ----
\newcommand{\figCorredor}{%
\begin{tikzpicture}
  \node[term] (L) at (0,0) {L};
  \node[st] (s1) at (1.6,0) {1};
  \node[st] (s2) at (3.2,0) {2};
  \node[st] (s3) at (4.8,0) {3};
  \node[term] (R) at (6.4,0) {R};
  \node[font=\footnotesize] at (0,-0.75) {entrar: $+2$};
  \node[font=\footnotesize] at (6.4,-0.75) {entrar: $+10$};
  \node[font=\footnotesize] at (3.2,-0.75) {cualquier otro movimiento: $-1$};
\end{tikzpicture}}

% ---- MDP del automóvil (clase) ----
\newcommand{\tabAuto}{%
{\small
\begin{tabular}{cccc}
\toprule
Estado & Acción & Sucesor (prob.) & Recompensa\\\midrule
Cool ($C$) & Slow & $C$ (1) & $+1$\\
Cool ($C$) & Fast & $C$ (0.5), $W$ (0.5) & $+2$\\
Warm ($W$) & Slow & $C$ (0.5), $W$ (0.5) & $+1$\\
Warm ($W$) & Fast & $O$ (1, terminal) & $-10$\\\bottomrule
\end{tabular}}}

\fancyhead[L]{\small\textcolor{gris}{Inteligencia Artificial 2027-I · Guía para el Examen 2}}
\fancyhead[R]{\small\textcolor{verde}{Respuestas}}
\fancyfoot[C]{\small\thepage}

\begin{document}

\begin{titlepage}
\centering
\vspace*{1.5cm}
{\color{verde}\rule{\linewidth}{1.2pt}}\\[0.6cm]
{\Huge\bfseries\color{verde} Guía para el Examen 2}\\[0.3cm]
{\LARGE Respuestas y soluciones desarrolladas}\\[0.2cm]
{\large Inteligencia Artificial · PCIC, UNAM · 2027-I}\\[0.4cm]
{\color{verde}\rule{\linewidth}{1.2pt}}\\[1.2cm]
\begin{minipage}{0.84\linewidth}\large
Este documento acompaña a la \emph{Guía de estudio para el Examen 2}. La numeración coincide
(A1–A9, B1–B34, C1–C13, D1–D2). En los ejercicios de procedimiento se muestran las tablas y trazas
completas, con el formato que conviene reproducir en el examen.

\vspace{0.5cm}
Todas las cifras (inferencia, muestreo, EM, regresión, $k$-means, iteración de valores y de
políticas, Q-learning, retropropagación y atención) se verificaron con un programa independiente.
Los valores se redondean a 4 cifras; los intermedios se arrastran con más precisión.

\vspace{0.5cm}
\textbf{Consejo:} usa estas soluciones para \emph{corregir}, no para leer. Si tu número no coincide,
busca el primer renglón de la tabla en el que te separas.
\end{minipage}
\vfill
\end{titlepage}

\tableofcontents
\newpage

% ============================================================
\section{Parte A — Conceptos}

\begin{solucion}{A1 · Probabilidad e incertidumbre}
\textbf{1.} \emph{Pereza}: enumerar todas las excepciones de una regla («dolor $\Rightarrow$ caries, salvo …») es
inviable. \emph{Ignorancia}: teórica (no hay teoría completa del dominio) y práctica (no se observa todo). La
probabilidad resume ambas en un grado de creencia, y la lógica pura no puede decidir entre acciones cuyo éxito no es seguro.
\textbf{2.} $P(\mathrm{caries})$ es marginal (suma 1 sobre los valores de Caries); $P(\mathrm{caries}\mid\mathrm{dolor})$ es condicional
(suma 1 sobre Caries \emph{para cada} valor fijo de Dolor); $P(\mathrm{caries},\mathrm{dolor})$ es conjunta (suma 1 sobre todas las parejas).
\textbf{3.} $3^8=6\,561$ filas y $6\,560$ parámetros libres.
\textbf{4.} Falso. En el colisionador $X\to Z\leftarrow Y$ (dos monedas y $Z$ = «salieron iguales»), $X\indep Y$ pero, dado $Z$,
conocer $X$ determina $Y$. También es falso el recíproco: en una cadena $X\to Z\to Y$, $X\indep Y\mid Z$ pero no $X\indep Y$.
\textbf{5.} La regla de la cadena factoriza \emph{cualquier} conjunta sin supuestos: $P(\mathbf x)=\prod_iP(x_i\mid x_{<i})$.
La red agrega el supuesto $P(x_i\mid x_{<i})=P(x_i\mid\mathrm{pa}(X_i))$, que reduce el número de parámetros.
\textbf{6.} Un modelo probabilístico (creencias $P(s'\mid a,\mathbf e)$) y una función de utilidad sobre resultados.
\end{solucion}

\begin{solucion}{A2 · Redes bayesianas e independencia}
\textbf{1.} La ausencia de arista \emph{no} garantiza independencia: puede haber caminos indirectos; lo que la red
garantiza son las independencias que se leen por d-separación. Una arista presente permite dependencia
pero no la obliga (los parámetros podrían hacerla nula).
\textbf{2.} Cada nodo es condicionalmente independiente de sus no descendientes dados sus padres. Manto de Markov:
padres, hijos y otros padres de los hijos; dado el manto, el nodo es independiente de todo lo demás.
\textbf{3.} El pasto está mojado: puede ser lluvia o aspersor. Si me dicen que el aspersor estuvo prendido, la
probabilidad de lluvia \emph{baja}: una causa ya explica el efecto (ver B6 y C1(d)).
\textbf{4.} Cadena y causa común: observar el nodo central \emph{bloquea}. Efecto común: observar el nodo central
(o un descendiente) \emph{activa}.
\textbf{5.} No necesariamente. Varias estructuras son equivalentes (mismas independencias: $X\to Y\to Z$ y
$X\leftarrow Y\leftarrow Z$) y puede haber variables no observadas; la causalidad requiere supuestos o intervenciones adicionales.
\textbf{6.} Si se agrega primero un efecto, sus causas deben depender de él y entre sí. Con orden $M,J,A,B,E$, $J$ necesita a $M$,
$A$ necesita $J,M$, $B$ necesita $A$, y $E$ necesita $A,B$: 13 parámetros en vez de 10. El orden causal suele dar la red más compacta.
\end{solucion}

\begin{solucion}{A3 · Inferencia exacta}
\textbf{1.} Consulta $Q$: la variable cuya distribución se busca; evidencia $\mathbf e$: valores observados; ocultas $\mathbf h$: el
resto, que se suman. $\alpha=1/P(\mathbf e)$ normaliza los pesos $P(q,\mathbf e)$ para que sumen 1.
\textbf{2.} En la enumeración, los productos de los factores que no dependen de la variable interior se recalculan para cada
valor de las exteriores. Eliminar guarda la suma parcial como un factor y la reutiliza (programación dinámica).
\textbf{3.} \emph{Restringir}: quedarse con las entradas compatibles con la evidencia. \emph{Unir}: multiplicar entradas
compatibles; el alcance es la unión de alcances. \emph{Eliminar}: sumar sobre los valores de una variable; el alcance pierde esa variable.
\textbf{4.} Suma y producto son conmutativos y distributivos, así que todo orden da el mismo número; cambia el tamaño de los
factores intermedios (B9).
\textbf{5.} Un poliárbol tiene a lo sumo un camino sin dirección entre dos nodos; en él la eliminación es lineal en el tamaño de
la red. En general la inferencia exacta es NP-difícil (\#P-difícil al contar).
\textbf{6.} Una oculta cuyos descendientes no están en la consulta ni en la evidencia: $\sum_hP(h\mid\mathrm{pa})=1$, de modo que
sus factores desaparecen. Se detecta podando hojas que no son consulta ni evidencia, repetidamente.
\end{solucion}

\begin{solucion}{A4 · Muestreo y aprendizaje de parámetros}
\textbf{1.} Se acepta una fracción $\approx P(\mathbf e)$ de las muestras: de $N$ se esperan $NP(\mathbf e)$ útiles. Con $P(\mathbf e)=10^{-3}$
casi todo se desperdicia, y el número crece exponencialmente con la cantidad de evidencia.
\textbf{2.} Las muestras se generan con la evidencia fija, no de $P(\mathbf z\mid\mathbf e)$; el peso $\prod P(e_i\mid\mathrm{pa})$ corrige esa
diferencia. Si la evidencia está abajo, no influye en el sorteo de sus ancestros: la mayoría de las muestras recibe peso diminuto
y pocas dominan (tamaño efectivo pequeño).
\textbf{3.} Solo el manto de Markov de la variable. Cada estado se obtiene modificando el anterior en una variable, por lo
que estados consecutivos se parecen; las primeras muestras dependen del estado inicial y se descartan (calentamiento).
\textbf{4.} $P(D\mid\theta)$ (o su logaritmo). Con datos completos el log se separa en una suma de términos, uno por fila de cada CPT,
cada uno maximizado por la frecuencia relativa.
\textbf{5.} Evita probabilidades cero por muestras pequeñas (que anulan productos y dan certeza falsa). Con $\lambda\to\infty$ cada fila
tiende a la uniforme (sesgo máximo).
\textbf{6.} E: con $\theta$ actual, calcula la posterior de las variables ocultas (conteos fraccionarios). M: reestima $\theta$ como MLE con esos
conteos. Garantiza que la verosimilitud observada no decrece; no garantiza el máximo global (depende de la inicialización).
\textbf{7.} Agregar padres nunca baja la verosimilitud de entrenamiento (la red más densa contiene a la más simple): se sobreajusta.
Hay que penalizar la complejidad o validar en datos no usados.
\end{solucion}

\begin{solucion}{A5 · Aprendizaje supervisado}
\textbf{1.} Clasificación predice una clase; regresión, una cantidad. «Lineal» significa lineal en los parámetros: $\hat y=\phi(x)^\mathsf T\theta$ con $\phi(x)=(1,x,x^2,\dots)$.
\textbf{2.} Si $X$ tiene rango columna completo. Si no (columnas dependientes o $n<d$), se usa pseudoinversa, regularización (ridge) o se eliminan columnas.
\textbf{3.} Subajuste: ambos errores altos (modelo demasiado simple). Sobreajuste: entrenamiento bajo y prueba alta (modelo que sigue el ruido).
\textbf{4.} La etiqueta es Bernoulli; maximizar su verosimilitud equivale a minimizar entropía cruzada. Esa pérdida es convexa en $\theta$
(hessiana $X^\mathsf T\mathrm{diag}(\pi(1-\pi))X\succeq0$); el error cuadrático con sigmoide no lo es y da gradientes pequeños.
\textbf{5.} Su frontera es una recta y XOR no es linealmente separable (B18). La capa oculta transforma el espacio en
características donde las clases sí se separan.
\textbf{6.} Batch usa todos los datos (gradiente exacto, caro); en línea uno (barato y ruidoso); mini-batch un subconjunto (compromiso, aprovecha GPU).
\textbf{7.} Una capa oculta suficientemente ancha puede aproximar cualquier función continua en un compacto. No dice cuántas
unidades, ni que el entrenamiento las encuentre, ni que generalice.
\textbf{8.} Una distribución predictiva: media y varianza en cada punto; la varianza crece lejos de los datos (útil para aprendizaje activo).
\end{solucion}

\begin{solucion}{A6 · Aprendizaje no supervisado}
\textbf{1.} La distancia suma diferencias en unidades distintas; la característica de mayor rango domina (C7).
\textbf{2.} Minimiza $J=\sum_j\sum_{x\in C_j}\|x-\mu_j\|^2$. La asignación y la actualización no aumentan $J$ y hay un número finito de
particiones, así que termina; solo en un óptimo local.
\textbf{3.} $0.8$: bien agrupado; $0$: en la frontera; $-0.3$: más cerca, en promedio, de otro grupo (probablemente mal asignado).
\textbf{4.} Simple: cadenas alargadas (une por el par más cercano). Completo: grupos compactos de diámetro pequeño.
\textbf{5.} Con grupos de forma arbitraria (anillos, medias lunas), con ruido o cuando no se conoce $K$.
\textbf{6.} Supervisado: usa etiquetas. $k$ = número de vecinos que votan; $K$ = número de grupos de $k$-means.
\end{solucion}

\begin{solucion}{A7 · MDP y aprendizaje por refuerzo}
\textbf{1.} $S$, $A(s)$, $T(s,a,s')=P(s'\mid s,a)$, $R(s,a,s')$, $\gamma$, estado inicial y terminales. Markov: el futuro depende solo del estado y la acción actuales.
\textbf{2.} Para que el retorno sea finito en horizontes infinitos y para preferir recompensas tempranas: $|G_t|\le R_{\max}/(1-\gamma)$.
\textbf{3.} $V^\pi(s)$: retorno esperado siguiendo $\pi$; $V^*$: el máximo sobre políticas; $Q^\pi(s,a)$, $Q^*(s,a)$: igual pero tomando primero $a$.
La política mapea estados a acciones. Con $Q^*$, $\pi^*(s)=\argmax_aQ^*(s,a)$ sin necesidad de $T$ ni $R$.
\textbf{4.} El operador de Bellman es una contracción en norma máximo: $\|BV-BU\|_\infty\le\gamma\|V-U\|_\infty$; tiene un único punto fijo y la iteración converge geométricamente.
\textbf{5.} Valores: un respaldo barato con $\max$ por iteración, muchas iteraciones. Políticas: evaluación exacta (sistema lineal, cara) y
mejora; pocas iteraciones y termina exactamente cuando la política no cambia.
\textbf{6.} Planificación: $T,R$ conocidos, se calcula sin actuar. RL: se aprende de experiencia; basado en modelo estima $T,R$ y planifica;
libre de modelo estima $V$ o $Q$ directamente.
\textbf{7.} MC necesita episodios completos; insesgado, varianza alta. TD actualiza en cada paso con $r+\gamma V(s')$ (\emph{bootstrap});
sesgado mientras $V$ sea inexacto, menor varianza, sirve en tareas continuas.
\textbf{8.} Q-learning usa $\max_{a'}Q(s',a')$: aprende el valor de la política voraz aunque se comporte explorando. SARSA usa la acción
realmente tomada: aprende el valor de la política que ejecuta.
\textbf{9.} Sin explorar, una mala estimación inicial nunca se corrige. Convergencia tabular: todos los pares $(s,a)$ visitados
infinitas veces y tasas con $\sum\alpha=\infty$, $\sum\alpha^2<\infty$.
\textbf{10.} Generalizar a estados no visitados con pocos parámetros; se pierde la garantía de convergencia y la capacidad de distinguir estados
con las mismas características. DQN: red neuronal como $\hat Q$, \emph{replay} (rompe correlaciones, reutiliza datos) y red objetivo congelada (estabiliza el objetivo).
\end{solucion}

\begin{solucion}{A8 · Aprendizaje profundo}
\textbf{1.} $W_2(W_1x+b_1)+b_2=(W_2W_1)x+(W_2b_1+b_2)$: sigue siendo afín.
\textbf{2.} La saturación de la sigmoide (derivada $\le1/4$, casi cero en los extremos) que desvanece gradientes. Neurona muerta: ReLU con
$z<0$ para todas las entradas relevantes; su gradiente es cero y no se recupera.
\textbf{3.} El gradiente de la pérdida respecto de todos los parámetros, aplicando la regla de la cadena desde la salida hacia la entrada y
reutilizando los $\boldsymbol\delta$; cuesta del orden de un forward.
\textbf{4.} Un vector aprendido por token (fila de $E\in\R^{V\times d}$). Coseno: ángulo entre vectores, $\frac{\mathbf u^\mathsf T\mathbf v}{\|\mathbf u\|\|\mathbf v\|}$; no es una probabilidad.
\textbf{5.} La consulta pregunta, las claves determinan la compatibilidad y los valores son lo que se mezcla. $\sqrt{d_k}$ evita puntajes
grandes que saturan softmax. La máscara impide ver tokens futuros en entrenamiento autorregresivo.
\textbf{6.} Preentrenamiento: predecir el siguiente token en mucho texto (autosupervisado). Instrucciones: ajuste supervisado con pares
instrucción–respuesta. RLHF: un modelo de recompensa entrenado con preferencias humanas guía un ajuste por RL.
\textbf{7.} $\tau<1$ concentra la distribución (más determinista); $\tau>1$ la aplana (más diversa); $\tau\to0$ es voraz.
\textbf{8.} Error esperado $=$ sesgo$^2$ $+$ varianza $+$ ruido. Modelos simples: sesgo alto; complejos: varianza alta; se elige con validación.
\end{solucion}

\begin{solucion}{A9 · IA en el contexto social}
\textbf{1.} Muestreo no representativo, etiquetas históricas sesgadas, medición (registrar reportes y no hechos), variable objetivo sustituta,
variables \emph{proxy}, y retroalimentación del despliegue.
\textbf{2.} Otras variables (código postal, historial) lo predicen, y las etiquetas pueden reflejar discriminación pasada.
\textbf{3.} Paridad: misma tasa de aprobación. Oportunidad: misma TPR. \emph{Equalized odds}: misma TPR y FPR. Con tasas base distintas,
igualar errores y mantener la calibración (PPV) a la vez es imposible salvo con un clasificador perfecto.
\textbf{4.} Texto en datos externos (correo, página, archivo) que el agente interpreta como orden. Controles: tratar contenido externo como
datos, mínimo privilegio, confirmación humana para acciones irreversibles, aislar herramientas y registrar acciones.
\textbf{5.} Inaceptable (prohibido: puntuación social estatal); alto riesgo (crédito, empleo, educación: obligaciones de gestión de riesgos,
datos, documentación, supervisión humana); riesgo limitado (chatbots, \emph{deepfakes}: transparencia); mínimo (filtros de spam, videojuegos).
\textbf{6.} W mide rapidez de consumo; kWh energía acumulada ($E=Pt$). Rebote: una eficiencia por uso que aumenta tanto el número de usos
que el consumo total crece (B34).
\textbf{7.} Exposición: tareas susceptibles de ser afectadas; sustitución depende de complementariedad y organización. Deuda cognitiva:
terminar tareas con ayuda sin construir la capacidad de resolverlas solo.
\textbf{8.} Delegación (qué delegar), descripción (comunicar bien la tarea), discernimiento (evaluar el resultado), diligencia (hacerse responsable).
\end{solucion}

% ============================================================
\newpage
\section{Parte B — Procedimientos}

\subsection{Probabilidad}

\begin{solucion}{B1 · Conjunta, marginales e independencia}
\textbf{(a)} $P(+t)=0.216+0.024+0.014+0.126=0.38$; $P(+p)=0.216+0.054+0.014+0.056=0.34$;
$P(+t,+p)=0.216+0.014=0.23$. Como $P(+t)P(+p)=0.1292\ne0.23$, \textbf{no} son independientes (el tráfico sugiere lluvia y la lluvia, paraguas).

\textbf{(b)} $P(+l\mid+t)=\dfrac{0.216+0.024}{0.38}=0.6316$; \quad $P(+l\mid+t,+p)=\dfrac{0.216}{0.23}=0.9391$. Dos indicios de la misma causa refuerzan la creencia.

\textbf{(c)} Con $P(+l)=0.3$ y $P(-l)=0.7$:
\begin{center}\small
\begin{tabular}{c|ccc}
 & $P(+t,+p\mid L)$ & $P(+t\mid L)$ & $P(+p\mid L)$\\\hline
$+l$ & $0.216/0.3=0.72$ & $0.24/0.3=0.8$ & $0.27/0.3=0.9$\\
$-l$ & $0.014/0.7=0.02$ & $0.14/0.7=0.2$ & $0.07/0.7=0.1$
\end{tabular}
\end{center}
$0.8\cdot0.9=0.72$ y $0.2\cdot0.1=0.02$; las otras tres combinaciones de cada estrato se siguen por complemento (p.\,ej.\ $P(+t,-p\mid+l)=0.024/0.3=0.08=0.8\cdot0.1$).
Por tanto $T\indep P\mid L$.

\textbf{(d)} Causa común $T\leftarrow L\to P$: $1+2+2=5$ parámetros frente a $2^3-1=7$.
\end{solucion}

\begin{solucion}{B2 · Bayes con dos pruebas}
\textbf{(a)} $P(+)=0.02(0.95)+0.98(0.08)=0.019+0.0784=0.0974$; $P(\text{enf}\mid+)=0.019/0.0974=\mathbf{0.1951}$.
A pesar de la buena sensibilidad, la mayoría de los positivos son falsos porque la enfermedad es rara.

\textbf{(b)} Conjunta: $\dfrac{0.02(0.95)^2}{0.02(0.95)^2+0.98(0.08)^2}=\dfrac{0.01805}{0.01805+0.006272}=\mathbf{0.7421}$.
Secuencial: usar $0.1951$ como nueva previa: $\dfrac{0.1951(0.95)}{0.1951(0.95)+0.8049(0.08)}=0.7421$. Coinciden porque las pruebas son
condicionalmente independientes dada la enfermedad.

\textbf{(c)} $\dfrac{0.02(0.95)(0.05)}{0.02(0.95)(0.05)+0.98(0.08)(0.92)}=\dfrac{0.00095}{0.073078}=\mathbf{0.0130}$.

\textbf{(d)} $\text{Prueba}_1\leftarrow\text{Enf}\to\text{Prueba}_2$ (causa común).
\end{solucion}

\begin{solucion}{B3 · Factorización y conteo}
\textbf{(a)} $P(A)P(B)P(C\mid A,B)P(D\mid C)P(E\mid C)P(F\mid E)$.
\textbf{(b)} $A$: $1$; $B$: $2$; $C$: $(2-1)\cdot2\cdot3=6$; $D$: $(4-1)\cdot2=6$; $E$: $(2-1)\cdot2=2$; $F$: $(3-1)\cdot2=4$. Total $\mathbf{21}$.
\textbf{(c)} Conjunta: $2\cdot3\cdot2\cdot4\cdot2\cdot3-1=\mathbf{287}$.
\textbf{(d)} $F$ pasa a tener padres $E,A$: $(3-1)\cdot2\cdot2=8$ en vez de 4: \textbf{4 parámetros más}.
\end{solucion}

\subsection{Redes bayesianas e independencia}

\begin{solucion}{B4 · Red del estudiante}
\textbf{(a)} $P(D)P(I)P(C\mid D,I)P(E\mid I)P(R\mid C)$: $1+1+4+2+2=10$ parámetros frente a $2^5-1=31$.

\textbf{(b)} $P(-d)P(+i)P(+c\mid-d,+i)P(-e\mid+i)P(+r\mid+c)=0.6(0.3)(0.9)(0.2)(0.9)=\mathbf{0.02916}$.

\textbf{(c)} $P(+c)=0.4(0.3)(0.7)+0.4(0.7)(0.1)+0.6(0.3)(0.9)+0.6(0.7)(0.4)=0.084+0.028+0.162+0.168=\mathbf{0.442}$.

\textbf{(d)} $D\indep I$: sí (único camino $D\to C\leftarrow I$, colisionador sin evidencia).
$D\indep I\mid R$: no garantizada ($R$ es descendiente de $C$ y abre el colisionador).
$D\indep E$: sí ($D\to C\leftarrow I\to E$ bloqueado en $C$).
$R\indep E\mid C$: sí ($R\leftarrow C\leftarrow I\to E$, cadena bloqueada en $C$ observado).
\end{solucion}

\begin{solucion}{B5 · d-separación}
Notación: \textcolor{verde}{A} = tripleta activa, \textcolor{rojo}{B} = bloqueada.
\begin{center}\small
\renewcommand{\arraystretch}{1.25}
\begin{tabularx}{\linewidth}{l X c}
\toprule
Consulta & Caminos y tripletas & ¿Garantizada?\\\midrule
(a) $A\indep B$ & $A\to C\leftarrow B$: colisionador $C$ sin evidencia ni descendientes observados \textcolor{rojo}{B}. $A\to C\to E\to G\leftarrow F\leftarrow D\leftarrow B$: colisionador $G$ \textcolor{rojo}{B}. & \textbf{Sí}\\
(b) $A\indep B\mid H$ & $A\to C\leftarrow B$: $H$ es descendiente de $C$ ($C\to E\to G\to H$) $\Rightarrow$ \textcolor{verde}{A}. & \textbf{No}\\
(c) $A\indep D$ & $A\to C\leftarrow B\to D$: $C$ \textcolor{rojo}{B}. $A\to C\to E\to G\leftarrow F\leftarrow D$: $G$ \textcolor{rojo}{B}. & \textbf{Sí}\\
(d) $A\indep D\mid G$ & $A\to C\leftarrow B\to D$: $C$ abierto por su descendiente $G$; $B$ (bifurcación) no observado \textcolor{verde}{A}. & \textbf{No}\\
(e) $A\indep D\mid G,B$ & Primer camino bloqueado en $B$. Pero $A\to C\to E\to G\leftarrow F\leftarrow D$: $C,E$ cadenas no observadas, $G$ colisionador observado, $F$ no observado \textcolor{verde}{A}. & \textbf{No}\\
(f) $C\indep F\mid B$ & $C\leftarrow B\to D\to F$: $B$ observado \textcolor{rojo}{B}. $C\to E\to G\leftarrow F$: $G$ y $H$ no observados \textcolor{rojo}{B}. & \textbf{Sí}\\
(g) $E\indep F\mid G,D$ & $E\to G\leftarrow F$: colisionador observado \textcolor{verde}{A}. (El camino por $C,B,D$ se bloquea en $D$.) & \textbf{No}\\
(h) $A\indep H\mid E$ & $A\to C\to E\to G\to H$: $E$ observado \textcolor{rojo}{B}. Pero $A\to C\leftarrow B\to D\to F\to G\to H$: $C$ abierto porque su descendiente $E$ está observado; el resto no observado \textcolor{verde}{A}. & \textbf{No}\\
(i) $A\indep F\mid C$ & $A\to C\leftarrow B\to D\to F$: $C$ colisionador observado \textcolor{verde}{A}. & \textbf{No}\\
(j) $B\indep G\mid E,F$ & $B\to C\to E\to G$: $E$ \textcolor{rojo}{B}. $B\to D\to F\to G$: $F$ \textcolor{rojo}{B}. $B\to C\leftarrow A$ no llega a $G$. & \textbf{Sí}\\\bottomrule
\end{tabularx}
\end{center}
\textbf{(k)} $\mathrm{MB}(C)=\{A,B,E\}$; $\mathrm{MB}(D)=\{B,F\}$; $\mathrm{MB}(F)=\{D,G,E\}$ ($E$ es otro padre de $G$); $\mathrm{MB}(G)=\{E,F,H\}$.
\begin{nota}
Las trampas son (e) y (h): observar algo puede bloquear un camino y, a la vez, abrir otro. Siempre revisa \emph{todos} los caminos.
\end{nota}
\end{solucion}

\begin{solucion}{B6 · Alarma: explicar a la otra causa}
Pesos $P(b,e,+a)=P(b)P(e)P(+a\mid b,e)$:
\begin{center}\small
\begin{tabular}{cc|l}
$B$ & $E$ & $P(b,e,+a)$\\\hline
$+b$ & $+e$ & $0.001(0.002)(0.95)=0.0000019$\\
$+b$ & $-e$ & $0.001(0.998)(0.94)=0.00093812$\\
$-b$ & $+e$ & $0.999(0.002)(0.29)=0.00057942$\\
$-b$ & $-e$ & $0.999(0.998)(0.001)=0.000997002$
\end{tabular}
\end{center}
\textbf{(a)} $P(+a)=\mathbf{0.002516442}$.
\textbf{(b)} $P(+b\mid+a)=\frac{0.0000019+0.00093812}{0.002516442}=\mathbf{0.3736}$; \quad $P(+e\mid+a)=\frac{0.0000019+0.00057942}{0.002516442}=\mathbf{0.2310}$.
\textbf{(c)} $P(+e\mid+a,+b)=\frac{0.0000019}{0.00094002}=\mathbf{0.0020}$; \quad $P(+e\mid+a,-b)=\frac{0.00057942}{0.001576422}=\mathbf{0.3676}$.
Saber que hubo robo «explica» la alarma y la probabilidad de terremoto cae de 0.23 a casi su previa (0.002).
\textbf{(d)} $J$ y $M$ son descendientes de $A$ no observados ni consultados: $\sum_jP(j\mid a)=\sum_mP(m\mid a)=1$.
\end{solucion}

\subsection{Inferencia exacta}

\begin{solucion}{B7 · Enumeración en una cadena}
\textbf{(a)}
\begin{center}\small
\begin{tabular}{ccc|l@{\qquad}ccc|l}
$R$&$T$&$L$& & $R$&$T$&$L$&\\\hline
$+$&$+$&$+$& $0.1(0.8)(0.3)=0.024$ & $-$&$+$&$+$& $0.9(0.1)(0.3)=0.027$\\
$+$&$+$&$-$& $0.056$ & $-$&$+$&$-$& $0.063$\\
$+$&$-$&$+$& $0.1(0.2)(0.1)=0.002$ & $-$&$-$&$+$& $0.9(0.9)(0.1)=0.081$\\
$+$&$-$&$-$& $0.018$ & $-$&$-$&$-$& $0.729$
\end{tabular}
\end{center}
\textbf{(b)} $P(+l)=0.024+0.002+0.027+0.081=0.134$.
$P(+r\mid+l)=\frac{0.024+0.002}{0.134}=\mathbf{0.1940}$ (sube desde 0.1). $P(+t\mid+l)=\frac{0.024+0.027}{0.134}=\mathbf{0.3806}$.
\textbf{(c)} $P(+l\mid+r)=0.8(0.3)+0.2(0.1)=\mathbf{0.26}$. Al condicionar en la raíz, $\sum_tP(t\mid+r)P(L\mid t)$ ya es una distribución
(producto de CPT que suman 1): normalizar dividiría entre $P(+r)/P(+r)=1$.
\end{solucion}

\begin{solucion}{B8 · Eliminación de variables en la red del estudiante}
\textbf{(a)} $P(D)$, $P(I)$, $P(C\mid D,I)$, $P(E\mid I)$, $P(+r\mid C)$ (este último es un factor $f(C)=(0.9,0.2)$).
\textbf{(b)} $E$ es hoja sin evidencia: $\sum_eP(e\mid I)=1$.
\textbf{(c)} Eliminar $D$: $f_1(C,I)=\sum_dP(d)P(C\mid d,I)$:
\begin{center}\small
\begin{tabular}{cc|l}
$C$ & $I$ & $f_1$\\\hline
$+c$ & $+i$ & $0.4(0.7)+0.6(0.9)=0.82$\\
$-c$ & $+i$ & $0.4(0.3)+0.6(0.1)=0.18$\\
$+c$ & $-i$ & $0.4(0.1)+0.6(0.4)=0.28$\\
$-c$ & $-i$ & $0.4(0.9)+0.6(0.6)=0.72$
\end{tabular}
\end{center}
\textbf{(d)} Eliminar $C$: $f_2(I)=\sum_cf_1(c,I)P(+r\mid c)$: $f_2(+i)=0.82(0.9)+0.18(0.2)=0.774$; $f_2(-i)=0.28(0.9)+0.72(0.2)=0.396$.
\textbf{(e)} $P(+i)f_2(+i)=0.2322$, $P(-i)f_2(-i)=0.2772$; suma $=P(+r)=0.5094$. $P(+i\mid+r)=\mathbf{0.4558}$.
\end{solucion}

\begin{solucion}{B9 · Orden de eliminación}
Factores iniciales: $P(A)$, $P(B)$, $P(C\mid A,B)$, $P(D\mid C)$, $P(E\mid C)$, $f(D,E)=P(+f\mid D,E)$.
\begin{center}\small
\begin{tabular}{l l l l}
\toprule
Orden & Elimina & Alcance del producto & Factor nuevo\\\midrule
(i) & $B$ & $\{A,B,C\}$ (8) & $f_1(A,C)$\\
 & $D$ & $\{C,D,E\}$ (8) & $f_2(C,E)$\\
 & $E$ & $\{C,E\}$ (4) & $f_3(C)$\\
 & $C$ & $\{A,C\}$ (4) & $f_4(A)$\\\midrule
(ii) & $C$ & $\{A,B,C,D,E\}$ (\textbf{32}) & $f_1(A,B,D,E)$ (16)\\
 & $B$ & $\{A,B,D,E\}$ (16) & $f_2(A,D,E)$\\
 & $D$ & $\{A,D,E\}$ (8) & $f_3(A,E)$\\
 & $E$ & $\{A,E\}$ (4) & $f_4(A)$\\\midrule
(iii) & $D$ & $\{C,D,E\}$ (8) & $f_1(C,E)$\\
 & $E$ & $\{C,E\}$ (4) & $f_2(C)$\\
 & $C$ & $\{A,B,C\}$ (8) & $f_3(A,B)$\\
 & $B$ & $\{A,B\}$ (4) & $f_4(A)$\\\bottomrule
\end{tabular}
\end{center}
Al final se multiplica por $P(A)$ y se normaliza. Mayor factor: (i) 8, (ii) 32, (iii) 8. Se recomiendan (i) o (iii): eliminar primero
$C$, que conecta a todos, crea un factor sobre todos sus vecinos.
\end{solucion}

\subsection{Muestreo}

\begin{solucion}{B10 · Muestreo previo y rechazo}
Regla: $+x$ si $u<P(+x\mid\text{padres})$.
\begin{center}\small
\begin{tabular}{c|llll|c}
\# & $C$ ($P(+c)=0.5$) & $S$ & $R$ & $W$ & Mundo\\\hline
1 & $0.31\Rightarrow+c$ & $0.62\ge0.1\Rightarrow-s$ & $0.12<0.8\Rightarrow+r$ & $0.45<0.90\Rightarrow+w$ & $+c,-s,+r,+w$\\
2 & $0.77\Rightarrow-c$ & $0.08<0.5\Rightarrow+s$ & $0.91\ge0.2\Rightarrow-r$ & $0.50<0.90\Rightarrow+w$ & $-c,+s,-r,+w$\\
3 & $0.18\Rightarrow+c$ & $0.35\Rightarrow-s$ & $0.66<0.8\Rightarrow+r$ & $0.97\ge0.90\Rightarrow-w$ & $+c,-s,+r,-w$\\
4 & $0.52\Rightarrow-c$ & $0.44<0.5\Rightarrow+s$ & $0.15<0.2\Rightarrow+r$ & $0.03<0.99\Rightarrow+w$ & $-c,+s,+r,+w$\\
5 & $0.05\Rightarrow+c$ & $0.06<0.1\Rightarrow+s$ & $0.70<0.8\Rightarrow+r$ & $0.88<0.99\Rightarrow+w$ & $+c,+s,+r,+w$
\end{tabular}
\end{center}
\textbf{(a)} $\hat P(+w)=4/5=\mathbf{0.8}$ (exacto 0.65).
\textbf{(b)} Se aceptan 1, 2, 4, 5; con $+c$: 1 y 5. $\hat P(+c\mid+w)=2/4=\mathbf{0.5}$ (exacto 0.5746).
\textbf{(c)} No: ninguna muestra tiene $-r,-w$; la razón sería $0/0$, indefinida (no es cero).
\end{solucion}

\begin{solucion}{B11 · Verosimilitud con pesos}
\textbf{(a)}
\begin{center}\small
\begin{tabular}{c|cc|l}
\# & $C$ & $R$ & $w=P(+s\mid c)P(+w\mid+s,r)$\\\hline
1 & $+c$ ($0.21$) & $+r$ ($0.40<0.8$) & $0.1(0.99)=0.099$\\
2 & $-c$ & $+r$ ($0.10<0.2$) & $0.5(0.99)=0.495$\\
3 & $+c$ & $-r$ ($0.93$) & $0.1(0.90)=0.09$\\
4 & $-c$ & $-r$ ($0.75$) & $0.5(0.90)=0.45$\\
5 & $+c$ & $-r$ ($0.85$) & $0.1(0.90)=0.09$
\end{tabular}
\end{center}
\textbf{(b)} $\hat P(+c\mid+s,+w)=\frac{0.099+0.09+0.09}{1.224}=\frac{0.279}{1.224}=\mathbf{0.2279}$.
\textbf{(c)} $P(+c,+s,+w)=0.5(0.1)[0.8(0.99)+0.2(0.9)]=0.0486$; $P(-c,+s,+w)=0.5(0.5)[0.2(0.99)+0.8(0.9)]=0.2295$;
exacto $0.0486/0.2781=\mathbf{0.1748}$. Con 5 muestras la diferencia es esperable.
\textbf{(d)} $S$ es evidencia: no se sortea, se fija. Su probabilidad depende del $C$ muestreado, y el peso corrige que el
procedimiento no la tuvo en cuenta al generar $C$ (sin él, $C$ seguiría su previa 0.5 y no la posterior).
\end{solucion}

\begin{solucion}{B12 · Dos barridos de Gibbs}
\textbf{(a)} Con $+w$ fijo:
$P(C\mid s,r)\propto P(C)P(s\mid C)P(r\mid C)$;\quad
$P(S\mid c,r,+w)\propto P(S\mid c)P(+w\mid S,r)$;\quad
$P(R\mid c,s,+w)\propto P(R\mid c)P(+w\mid s,R)$.

\textbf{(b)}
\begin{center}\small
\begin{tabular}{c c l l c c}
\toprule
Paso & Var. & Peso $+$ & Peso $-$ & $P(+)$ & Nuevo estado $(C,S,R)$\\\midrule
1 & $C$ & $0.5(0.9)(0.2)=0.09$ & $0.5(0.5)(0.8)=0.2$ & 0.3103 & $u=0.60\Rightarrow(-c,-s,-r)$\\
2 & $S$ & $0.5(0.90)=0.45$ & $0.5(0.01)=0.005$ & 0.9890 & $u=0.20\Rightarrow(-c,+s,-r)$\\
3 & $R$ & $0.2(0.99)=0.198$ & $0.8(0.90)=0.72$ & 0.2157 & $u=0.30\Rightarrow(-c,+s,-r)$\\
4 & $C$ & $0.5(0.1)(0.2)=0.01$ & $0.5(0.5)(0.8)=0.2$ & 0.0476 & $u=0.90\Rightarrow(-c,+s,-r)$\\
5 & $S$ & $0.45$ & $0.005$ & 0.9890 & $u=0.55\Rightarrow(-c,+s,-r)$\\
6 & $R$ & $0.198$ & $0.72$ & 0.2157 & $u=0.40\Rightarrow(-c,+s,-r)$\\\bottomrule
\end{tabular}
\end{center}
\textbf{(c)} Ambos estados tras cada barrido tienen $-c$: $\hat P(+c\mid+w)=0$, lejos de 0.5746. Son pocas muestras, muy correlacionadas
y sin calentamiento; la cadena se quedó en una región ($+s$ explica $+w$). Hace falta correr mucho más.
\end{solucion}

\subsection{Aprendizaje de parámetros}

\begin{solucion}{B13 · MLE y Laplace con dos padres}
\textbf{(a)} Conteos: $S=5$, $N=4$, $Ll=3$. MLE: $(5/12,4/12,3/12)=(0.4167,0.3333,0.25)$. Laplace ($k=3$): $(6/15,5/15,4/15)=(0.4,0.3333,0.2667)$.
\textbf{(b), (c)} Filas $P(\text{sí}\mid\text{Clima},\text{Día})$ (Laplace con $k=2$: $\frac{N(\text{sí})+1}{N+2}$):
\begin{center}\small
\begin{tabular}{cc|ccc}
Clima & Día & $N$ (sí/total) & MLE & Laplace\\\hline
S & L & 2/3 & 0.6667 & 3/5 = 0.6\\
S & F & 2/2 & 1 & 3/4 = 0.75\\
N & L & 1/2 & 0.5 & 2/4 = 0.5\\
N & F & 1/2 & 0.5 & 2/4 = 0.5\\
Ll & L & 1/3 & 0.3333 & 2/5 = 0.4\\
Ll & F & 0/0 & \textbf{indefinida} & 1/2 = 0.5
\end{tabular}
\end{center}
La fila (Ll, F) nunca aparece; con Laplace queda uniforme. Nota que MLE da $P(\text{no}\mid S,F)=0$: certeza sacada de dos datos.
\textbf{(d)} Clima 2 + Día 1 + Bici $1\cdot3\cdot2=6$: \textbf{9} parámetros libres.
\end{solucion}

\begin{solucion}{B14 · Una iteración de EM}
\textbf{(a)} Paso E: $w_j=\frac{P(a)P(r_1\mid a)P(r_2\mid a)}{P(a)P(r_1\mid a)P(r_2\mid a)+P(b)P(r_1\mid b)P(r_2\mid b)}$.
\begin{center}\small
\begin{tabular}{c|ll|c}
Registro & $a$: $0.5\cdot\cdot$ & $b$: $0.5\cdot\cdot$ & $w_j$\\\hline
$(1,1)$ & $0.5(0.8)(0.7)=0.28$ & $0.5(0.3)(0.4)=0.06$ & 0.8235\\
$(1,0)$ & $0.5(0.8)(0.3)=0.12$ & $0.5(0.3)(0.6)=0.09$ & 0.5714\\
$(0,0)$ & $0.5(0.2)(0.3)=0.03$ & $0.5(0.7)(0.6)=0.21$ & 0.1250\\
$(1,1)$ & 0.28 & 0.06 & 0.8235\\
$(0,1)$ & $0.5(0.2)(0.7)=0.07$ & $0.5(0.7)(0.4)=0.14$ & 0.3333
\end{tabular}
\end{center}
\textbf{(b)} $\sum_jw_j=2.6768$, $\sum_j(1-w_j)=2.3232$.
$P(a)=2.6768/5=\mathbf{0.5354}$.
$P(R_1{=}1\mid a)=\frac{0.8235+0.5714+0.8235}{2.6768}=\mathbf{0.8288}$; \quad
$P(R_1{=}1\mid b)=\frac{0.1765+0.4286+0.1765}{2.3232}=\mathbf{0.3364}$.
$P(R_2{=}1\mid a)=\frac{0.8235+0.8235+0.3333}{2.6768}=\mathbf{0.7398}$; \quad
$P(R_2{=}1\mid b)=\frac{0.1765+0.1765+0.6667}{2.3232}=\mathbf{0.4389}$.
\textbf{(c)} Antes: $P(r_1,r_2)$ = suma de las dos columnas: $0.34,0.21,0.24,0.34,0.21$; $\sum\log=\mathbf{-6.7059}$.
Después: $\mathbf{-6.6613}$. Aumenta, como garantiza EM (la siguiente iteración da $-6.6611$: ya casi convergió).
\end{solucion}

\subsection{Aprendizaje supervisado}

\begin{solucion}{B15 · Regresión lineal y descenso de gradiente}
\textbf{(a)} $X=\begin{pmatrix}1&0\\1&1\\1&2\\1&3\end{pmatrix}$, $X^\mathsf TX=\begin{pmatrix}4&6\\6&14\end{pmatrix}$, $X^\mathsf T\mathbf y=\begin{pmatrix}16\\35\end{pmatrix}$, $\det=20$.
$\hat\theta=\frac1{20}\begin{pmatrix}14&-6\\-6&4\end{pmatrix}\begin{pmatrix}16\\35\end{pmatrix}=\begin{pmatrix}\mathbf{0.7}\\\mathbf{2.2}\end{pmatrix}$.

\textbf{(b)} $\hat y=(0.7,2.9,5.1,7.3)$; residuos $(0.3,0.1,-1.1,0.7)$; $J=0.09+0.01+1.21+0.49=\mathbf{1.8}$. $\hat y(5)=0.7+11=\mathbf{11.7}$.

\textbf{(c)} $\nabla J=2X^\mathsf T(X\theta-\mathbf y)$.
Paso 1: en $\theta=(0,0)$, $\nabla J=-2(16,35)=(-32,-70)$; $\theta_1=(0.32,0.70)$; $J_0=90$, $J_1=\mathbf{40.7176}$.
Paso 2: residuos $X\theta_1-\mathbf y=(-0.68,-1.98,-2.28,-5.58)$; $\nabla J=(-21.04,-46.56)$; $\theta_2=(0.5304,1.1656)$; $J_2=\mathbf{19.0000}$.
La pérdida baja y $\theta$ se acerca a $(0.7,2.2)$.
\end{solucion}

\begin{solucion}{B16 · Regresión logística}
\textbf{(a)} $z=X\theta=(1,-2,1)$; $\pi=(0.7311,0.1192,0.7311)$; pérdidas $-\ln0.7311=0.3133$, $-\ln(1-0.1192)=0.1269$, $0.3133$; $J=\mathbf{0.7535}$.
\textbf{(b)} $\boldsymbol\pi-\mathbf y=(-0.2689,0.1192,-0.2689)$.
$\nabla J=\big(\sum(\pi-y),\ \sum x_1(\pi-y),\ \sum x_2(\pi-y)\big)=\big(-0.4187,\ 2(-0.2689)+1(-0.2689),\ 1(-0.2689)+2(0.1192)\big)=(\mathbf{-0.4187},\mathbf{-0.8068},\mathbf{-0.0305})$.
\textbf{(c)} $\theta_1=\theta-0.5\nabla J=(0.2093,1.4034,-0.9847)$; $z=(2.0314,-1.7601,1.6128)$, $\pi=(0.8841,0.1468,0.8338)$; $J_1=\mathbf{0.4637}$.
\textbf{(d)} Antes: $x_1-x_2=0$, es decir $x_2=x_1$. Después: $0.2093+1.4034x_1-0.9847x_2=0$, es decir $x_2=0.2126+1.4252x_1$.
\end{solucion}

\begin{solucion}{B17 · Softmax}
\textbf{(a)} $e^{\mathbf z}=(7.3891,2.7183,0.3679)$, suma $10.4752$; $\mathbf p=(0.7054,0.2595,0.0351)$; $\ell=-\ln0.2595=\mathbf{1.3490}$.
\textbf{(b)} $\mathbf p-\mathbf t=(0.7054,-0.7405,0.0351)$: sube el logit correcto, bajan los demás.
\textbf{(c)} $\tau=2$: softmax de $(1,0.5,-0.5)$ $=(0.5465,0.3315,0.1220)$. $\tau=0.5$: softmax de $(4,2,-2)$ $=(0.8789,0.1189,0.0022)$.
Con $\tau\to0$ toda la masa va al máximo (voraz).
\textbf{(d)} $\frac{e^{z_j-m}}{\sum_ke^{z_k-m}}=\frac{e^{-m}e^{z_j}}{e^{-m}\sum_ke^{z_k}}=p_j$; evita desbordes numéricos.
\end{solucion}

\begin{solucion}{B18 · XOR}
\textbf{(a)} Supón salida 1 si $w_1x_1+w_2x_2+b>0$. Se requiere $b\le0$ (para $00$), $w_2+b>0$ (01), $w_1+b>0$ (10), $w_1+w_2+b\le0$ (11).
Sumando las del medio: $w_1+w_2+2b>0\Rightarrow w_1+w_2+b>-b\ge0$, contradicción con la última.
\textbf{(b)}
\begin{center}\small
\begin{tabular}{cc|cc|c}
$x_1$&$x_2$&$h_1$&$h_2$&$h_1-2h_2$\\\hline
0&0&0&0&0\\0&1&1&0&1\\1&0&1&0&1\\1&1&2&1&0
\end{tabular}
\end{center}
\textbf{(c)} Oculta: 4 pesos + 2 sesgos; salida: 2 pesos + 1 sesgo (0): \textbf{9}. Las dos unidades miden «al menos uno encendido» y
«ambos encendidos»; en el espacio $(h_1,h_2)$ los puntos $01$ y $10$ colapsan en $(1,0)$ y la salida es lineal.
\end{solucion}

\subsection{No supervisado y vecinos}

\begin{solucion}{B19 · Distancias y $k$-NN}
\textbf{(a)}
\begin{center}\small
\begin{tabular}{c|c|ccc}
Punto & $x-q$ & $L_1$ & $L_2$ & $L_\infty$\\\hline
$P_1(6,3)[+]$ & $(3,0)$ & 3 & 3.000 & 3\\
$P_2(5,5)[-]$ & $(2,2)$ & 4 & 2.828 & 2\\
$P_3(-1,4)[+]$ & $(-4,1)$ & 5 & 4.123 & 4\\
$P_4(1,-1)[-]$ & $(-2,-4)$ & 6 & 4.472 & 4
\end{tabular}
\end{center}
\textbf{(b)} $L_1$: 1-NN $P_1\Rightarrow\mathbf{+}$; 3-NN $\{P_1,P_2,P_3\}\Rightarrow\mathbf{+}$ (2 a 1).
$L_2$: 1-NN $P_2\Rightarrow\mathbf{-}$; 3-NN $\{P_2,P_1,P_3\}\Rightarrow\mathbf{+}$.
$L_\infty$: 1-NN $P_2\Rightarrow\mathbf{-}$; 3-NN: $P_2,P_1$ y empate $P_3/P_4$ a distancia 4: la regla de desempate decide ($+$ con $P_3$, $-$ con $P_4$).
La métrica cambia la respuesta.
\textbf{(c)} Difieren en tamaño y material: Hamming $=\mathbf2$.
\end{solucion}

\begin{solucion}{B20 · $k$-means}
\textbf{Iteración 1} ($\mu_1=(1,1)$, $\mu_2=(1,3)$), $d^2$:
\begin{center}\small
\begin{tabular}{c|cccccc}
 & $P_1$ & $P_2$ & $P_3$ & $P_4$ & $P_5$ & $P_6$\\\hline
$\mu_1$ & \textbf{0} & \textbf{1} & 4 & 25 & 41 & 41\\
$\mu_2$ & 4 & 5 & \textbf{0} & \textbf{17} & \textbf{29} & \textbf{25}
\end{tabular}
\end{center}
Grupos $\{P_1,P_2\}$, $\{P_3,P_4,P_5,P_6\}$; $J=72$. Actualización: $\mu_1=(1.5,1)$, $\mu_2=(4.25,4.5)$; $J=20.25$.

\textbf{Iteración 2}:
\begin{center}\small
\begin{tabular}{c|cccccc}
 & $P_1$ & $P_2$ & $P_3$ & $P_4$ & $P_5$ & $P_6$\\\hline
$\mu_1$ & \textbf{0.25} & \textbf{0.25} & \textbf{4.25} & 21.25 & 36.25 & 37.25\\
$\mu_2$ & 22.81 & 17.31 & 12.81 & \textbf{0.81} & \textbf{3.31} & \textbf{2.81}
\end{tabular}
\end{center}
$P_3$ cambia de grupo. $J=11.6875$. Actualización: $\mu_1=(4/3,5/3)\approx(1.333,1.667)$, $\mu_2=(16/3,5)\approx(5.333,5)$; $J=\mathbf{6.0}$.

\textbf{Iteración 3}: distancias de $P_1,P_2,P_3$ a $\mu_1$: $0.56,0.89,1.89$ (a $\mu_2$: $34.8,27.1,22.8$); de $P_4,P_5,P_6$ a $\mu_2$: $1.11,0.44,1.11$.
Sin cambios: \textbf{convergió} con $\{P_1,P_2,P_3\}$, $\{P_4,P_5,P_6\}$, $J=6$.
\end{solucion}

\begin{solucion}{B21 · Silueta}
$P_3=(1,3)$: $a=\frac{d(P_3,P_1)+d(P_3,P_2)}2=\frac{2+2.236}{2}=2.118$; $b=\frac{4.123+5.385+5}{3}=4.836$; $s=\frac{4.836-2.118}{4.836}=\mathbf{0.562}$.
$P_1=(1,1)$: $a=\frac{1+2}{2}=1.5$; $b=\frac{5+6.403+6.403}3=5.935$; $s=\mathbf{0.747}$. $P_1$ está mejor agrupado; $P_3$ es el más periférico de su grupo.
\end{solucion}

\begin{solucion}{B22 · Agrupamiento jerárquico}
\begin{center}\small
\begin{tabular}{c|ll|ll|ll}
\toprule
Paso & Simple & altura & Completo & altura & Promedio & altura\\\midrule
1 & $A+B$ & 2 & $A+B$ & 2 & $A+B$ & 2\\
2 & $D+E$ & 3 & $D+E$ & 3 & $D+E$ & 3\\
3 & $C+DE$ & $\min(4,5)=4$ & $C+DE$ & $\max(4,5)=5$ & $C+DE$ & $\frac{4+5}2=4.5$\\
4 & $AB+CDE$ & $\min=5$ & $AB+CDE$ & $\max=10$ & $AB+CDE$ & $\frac{47}6=7.83$\\\bottomrule
\end{tabular}
\end{center}
En el paso 3 se compara $C$ con $AB$ (simple 5, completo 6, promedio 5.5) frente a $C$ con $DE$; gana $DE$ en los tres.
En el paso 4 promedio: $(6+10+9+5+9+8)/6=47/6$. Corte en 4.2: simple $\{AB\},\{CDE\}$ (\textbf{2}); completo y promedio $\{AB\},\{C\},\{DE\}$ (\textbf{3}).
Mismo orden de uniones, distintas alturas: el corte cambia el número de grupos.
\end{solucion}

\subsection{Aprendizaje por refuerzo}

\begin{solucion}{B23 · Retornos, MC y TD(0)}
\textbf{(a)} Hacia atrás con $\gamma=0.5$. Ep.\,1: $G_3=4$ ($B$), $G_2=-1+0.5(4)=1$ ($A$), $G_1=2+0.5(1)=2.5$ ($B$), $G_0=1+0.5(2.5)=2.25$ ($A$).
Ep.\,2: $G_1=1$ ($A$), $G_0=3+0.5(1)=3.5$ ($B$).
\textbf{(b)} Primera visita: $V(A)=\frac{2.25+1}{2}=\mathbf{1.625}$, $V(B)=\frac{2.5+3.5}2=\mathbf{3.0}$.
Todas las visitas: $V(A)=\frac{2.25+1+1}3=\mathbf{1.4167}$, $V(B)=\frac{2.5+4+3.5}3=\mathbf{3.3333}$.
\textbf{(c)} $V(s)\leftarrow V(s)+0.5[r+0.5V(s')-V(s)]$:
\begin{center}\small
\begin{tabular}{c|c|c|cc}
Transición & Objetivo & Actualiza & $V(A)$ & $V(B)$\\\hline
$A\xrightarrow{1}B$ & $1+0.5(0)=1$ & $A$: $0\to0.5$ & 0.5 & 0\\
$B\xrightarrow{2}A$ & $2+0.5(0.5)=2.25$ & $B$: $0\to1.125$ & 0.5 & 1.125\\
$A\xrightarrow{-1}B$ & $-1+0.5(1.125)=-0.4375$ & $A$: $0.5\to0.03125$ & 0.0313 & 1.125\\
$B\xrightarrow{4}T$ & $4$ & $B$: $1.125\to2.5625$ & 0.0313 & 2.5625\\
$B\xrightarrow{3}A$ & $3+0.5(0.03125)=3.0156$ & $B$: $\to2.7891$ & 0.0313 & 2.7891\\
$A\xrightarrow{1}T$ & $1$ & $A$: $\to0.5156$ & 0.5156 & 2.7891
\end{tabular}
\end{center}
\end{solucion}

\begin{solucion}{B24 · Evaluación exacta de «siempre derecha»}
\textbf{(a)} $V_1=0.8(-1+0.9V_2)+0.2(2)=-0.4+0.72V_2$;\quad
$V_2=0.8(-1+0.9V_3)+0.2(-1+0.9V_1)=-1+0.72V_3+0.18V_1$;\quad
$V_3=0.8(10)+0.2(-1+0.9V_2)=7.8+0.18V_2$.
Sustituyendo $V_1$ y $V_3$ en $V_2$: $V_2=-1+0.72(7.8+0.18V_2)+0.18(-0.4+0.72V_2)=4.544+0.2592V_2$, así
$V_2=\mathbf{6.1339}$, $V_1=\mathbf{4.0164}$, $V_3=\mathbf{8.9041}$.

\textbf{(b)} $Q^\pi(s,a)=\sum T[R+0.9V^\pi(s')]$:
\begin{center}\small
\begin{tabular}{c|cc}
$s$ & izq & der\\\hline
1 & $0.8(2)+0.2(-1+0.9V_2)=2.5041$ & 4.0164\\
2 & $0.8(-1+0.9V_1)+0.2(-1+0.9V_3)=3.4946$ & 6.1339\\
3 & $0.8(-1+0.9V_2)+0.2(10)=5.6164$ & 8.9041
\end{tabular}
\end{center}
Ningún estado mejora: la política es voraz respecto de su propio valor, satisface Bellman de optimalidad y es \textbf{óptima}.
\end{solucion}

\begin{solucion}{B25 · Iteración de valores en el corredor}
\begin{center}\small
\begin{tabular}{c|cc|cc|cc|c}
\toprule
$k$ & $Q(1,\text{izq})$ & $Q(1,\text{der})$ & $Q(2,\text{izq})$ & $Q(2,\text{der})$ & $Q(3,\text{izq})$ & $Q(3,\text{der})$ & $V_k=(V_1,V_2,V_3)$\\\midrule
1 & 1.4 & $-0.4$ & $-1$ & $-1$ & 1.2 & 7.8 & $(1.4,-1,7.8)$\\
2 & 1.22 & $-1.12$ & 1.412 & 4.868 & 0.48 & 7.62 & $(1.22,4.868,7.62)$\\
3 & 2.2762 & 3.105 & 1.25 & 4.706 & 4.705 & 8.6762 & $(3.105,4.706,8.6762)$\\\bottomrule
\end{tabular}
\end{center}
Ejemplo ($k=2$): $Q(2,\text{der})=0.8(-1+0.9\cdot7.8)+0.2(-1+0.9\cdot1.4)=0.8(6.02)+0.2(0.26)=4.868$.
Políticas voraces: $k=1$: (izq, empate, der); $k=2$: (izq, der, der); $k=3$: (der, der, der), que ya es la óptima.
$V_k$ es el mejor retorno con $k$ pasos: desde 1 hacen falta 3 movimientos para entrar a $R$, así que con $k\le2$ el $+2$ de $L$ es lo mejor alcanzable.
\end{solucion}

\begin{solucion}{B26 · Q-learning en el automóvil}
\textbf{(a)} $Q\leftarrow Q+0.5[y-Q]$, $y=r+0.9\max_{a'}Q(s',a')$ ($O$ terminal: $y=r$).
\begin{center}\small
\begin{tabular}{c|c|c|c}
Transición & $y$ & Anterior & Nuevo\\\hline
$(C,F,+2,W)$ & $2+0.9\max(0,0)=2$ & 0 & $Q(C,F)=1$\\
$(W,S,+1,C)$ & $1+0.9\max(0,1)=1.9$ & 0 & $Q(W,S)=0.95$\\
$(C,F,+2,C)$ & $2+0.9(1)=2.9$ & 1 & $Q(C,F)=1.95$\\
$(C,S,+1,C)$ & $1+0.9(1.95)=2.755$ & 0 & $Q(C,S)=1.3775$\\
$(C,F,+2,W)$ & $2+0.9(0.95)=2.855$ & 1.95 & $Q(C,F)=2.4025$\\
$(W,F,-10,O)$ & $-10$ & 0 & $Q(W,F)=-5$
\end{tabular}
\end{center}
\textbf{(b)} $C$: Fast ($2.4025>1.3775$); $W$: Slow ($0.95>-5$).
\textbf{(c)} Coincide con la óptima (Fast en Cool, Slow en Warm), pero los valores están lejos de $Q^*$ ($V^*=(15.5,14.5)$ con $\gamma=0.9$).
Para converger: visitar todos los pares muchas veces, explorar y reducir $\alpha$ adecuadamente.
\end{solucion}

\begin{solucion}{B27 · $\varepsilon$-voraz, SARSA y Q-learning}
\textbf{(a)} $\mathcal G=\{a_2,a_3\}$, $\varepsilon/m=0.05$, $(1-\varepsilon)/|\mathcal G|=0.4$: $\mu=(\mathbf{0.05},\mathbf{0.45},\mathbf{0.45},\mathbf{0.05})$.
\textbf{(b)} Q-learning: $y=0+0.9\max(6,-2)=5.4$; $Q=2+0.5(3.4)=\mathbf{3.7}$. SARSA: $y=0+0.9(-2)=-1.8$; $Q=2+0.5(-3.8)=\mathbf{0.1}$.
Q-learning aprende el valor de la voraz (fuera de política); SARSA el de la política que explora (en política), por eso es más
conservador cerca de acciones peligrosas.
\end{solucion}

\begin{solucion}{B28 · Q-learning lineal}
\textbf{(a)} $\hat Q(s,a)=1+2-2=1$. En $s'$: $1+0-1=0$ y $1+4+0=5$; objetivo $y=3+0.9(5)=7.5$; $\delta=6.5$.
\textbf{(b)} $\mathbf w'=(1,2,-1)+0.1(6.5)(1,1,2)=(\mathbf{1.65},\mathbf{2.65},\mathbf{0.3})$; ahora $\hat Q(s,a)=4.9$.
\textbf{(c)} En $s'$ pasan a $1.95$ y $6.95$ (ambos $+1.95$): generalización por las características compartidas (sesgo y $\phi_1$).
Semigradiente: se deriva $\frac12(y-\hat Q)^2$ tratando $y$ como constante, aunque también depende de $\mathbf w$.
\end{solucion}

\subsection{Aprendizaje profundo y evaluación}

\begin{solucion}{B29 · Retropropagación 2–2–1}
\textbf{(a)} $\mathbf z_1=W_1\mathbf x+\mathbf b_1=(1-2+0,\ 0.5+1-1)=(-1,0.5)$; $\mathbf h=(0,0.5)$; $z_2=2(0)-1(0.5)+0.5=0$; $p=0.5$; $\ell=\ln2=\mathbf{0.6931}$.
\textbf{(b)} $\delta_2=p-y=-0.5$; $\partial\ell/\partial\mathbf w_2=\delta_2\mathbf h=(0,-0.25)$; $\partial\ell/\partial b_2=-0.5$.
$\boldsymbol\delta_1=(\mathbf w_2\delta_2)\odot\mathrm{ReLU}'(\mathbf z_1)=(2(-0.5)\cdot0,\ -1(-0.5)\cdot1)=(0,0.5)$.
$\partial\ell/\partial W_1=\boldsymbol\delta_1\mathbf x^\mathsf T=\begin{pmatrix}0&0\\0.5&1\end{pmatrix}$; $\partial\ell/\partial\mathbf b_1=(0,0.5)$.
\textbf{(c)} $\eta=0.5$: $W_1=\begin{pmatrix}1&-1\\0.25&0\end{pmatrix}$, $\mathbf b_1=(0,-1.25)$, $\mathbf w_2=(2,-0.875)$, $b_2=0.75$.
Nuevo forward: $\mathbf z_1=(-1,-1)$, $\mathbf h=(0,0)$, $z_2=0.75$, $p=0.6792$, $\ell=\mathbf{0.3869}$. La pérdida baja (aunque ahora ambas ocultas se apagaron
para este ejemplo y solo $b_2$ podrá seguir aprendiendo con él).
\textbf{(d)} La primera unidad oculta tiene $z=-1<0$: ReLU$'=0$ corta el gradiente hacia su fila de $W_1$ y su sesgo.
\end{solucion}

\begin{solucion}{B30 · Conteo y composición}
\textbf{(a)} $(8\cdot4+8)+(8\cdot8+8)+(3\cdot8+3)=40+72+27=\mathbf{139}$.
\textbf{(b)} $32\,000\cdot256=\mathbf{8\,192\,000}$.
\textbf{(c)} $W_2(W_1\mathbf x+\mathbf b_1)+\mathbf b_2=(W_2W_1)\mathbf x+(W_2\mathbf b_1+\mathbf b_2)=\widetilde W\mathbf x+\widetilde{\mathbf b}$.
\textbf{(d)} $(3\cdot2+3)+(2\cdot3+2)=\mathbf{17}$. Cada unidad ReLU divide el plano con una recta; 3 rectas en posición general forman a lo sumo $1+3+3=7$ regiones, en cada una la red es afín.
\end{solucion}

\begin{solucion}{B31 · Atención causal}
\textbf{(a)} $S=\frac1{\sqrt2}\begin{pmatrix}1&1&0\\0&1&1\\1&2&1\end{pmatrix}=\begin{pmatrix}0.7071&0.7071&0\\0&0.7071&0.7071\\0.7071&1.4142&0.7071\end{pmatrix}$.
\textbf{(b)} Con máscara, fila 1: $(1,0,0)$; fila 2: softmax$(0,0.7071)=(0.3302,0.6698,0)$; fila 3: softmax$(0.7071,1.4142,0.7071)=(0.2483,0.5035,0.2483)$.
\textbf{(c)} $\mathbf o_1=(1,0)$; $\mathbf o_2=0.3302(1,0)+0.6698(0,1)=(0.3302,0.6698)$; $\mathbf o_3=0.2483(1,0)+0.5035(0,1)+0.2483(1,1)=(0.4966,0.7517)$.
\textbf{(d)} Sin máscara la fila 2 sería softmax$(0,0.7071,0.7071)=(0.1978,0.4011,0.4011)$ y $\mathbf o_2=(0.5989,0.8022)$: la posición 2 «vería el futuro».
\end{solucion}

\begin{solucion}{B32 · Matriz de confusión y umbral}
\begin{center}\small
\begin{tabular}{c|cccc|ccccc}
\toprule
Umbral & $TP$ & $FP$ & $FN$ & $TN$ & Exactitud & Precisión & Recall & Especif. & $F_1$\\\midrule
0.50 & 3 & 2 & 2 & 3 & 0.6 & 0.6 & 0.6 & 0.6 & 0.6\\
0.35 & 4 & 3 & 1 & 2 & 0.6 & 0.5714 & 0.8 & 0.4 & 0.6667\\
0.75 & 2 & 1 & 3 & 4 & 0.6 & 0.6667 & 0.4 & 0.8 & 0.5\\\bottomrule
\end{tabular}
\end{center}
La exactitud no distingue los tres umbrales. Enfermedad grave: un falso negativo es muy costoso $\Rightarrow$ umbral bajo (0.35, recall 0.8).
Spam: perder un correo legítimo es costoso $\Rightarrow$ umbral alto (0.75, precisión y especificidad mayores).
\end{solucion}

\subsection{IA en el contexto social}

\begin{solucion}{B33 · Métricas de equidad}
\textbf{(a)}
\begin{center}\small
\begin{tabular}{c|cccccc}
Grupo & Tasa base & Aprobación & TPR & FPR & PPV & Exactitud\\\hline
A & $50/100=0.5$ & $50/100=0.5$ & $40/50=0.8$ & $10/50=0.2$ & $40/50=0.8$ & 0.8\\
B & $30/100=0.3$ & $20/100=0.2$ & $15/30=0.5$ & $5/70=0.0714$ & $15/20=0.75$ & 0.8
\end{tabular}
\end{center}
\textbf{(b)} Paridad: no ($0.5$ vs $0.2$). Igualdad de oportunidades: no ($0.8$ vs $0.5$: la mitad de los buenos pagadores de B son rechazados).
\emph{Equalized odds}: no (además FPR $0.2$ vs $0.071$).
\textbf{(c)} No. La exactitud igual oculta \emph{quién} sufre los errores: en B dominan los falsos negativos (personas que pagarían y no reciben
crédito). Hay que reportar métricas desagregadas y decidir qué error importa.
\end{solucion}

\begin{solucion}{B34 · Energía y rebote}
\textbf{(a)} $10\,\mathrm W\times100\,\mathrm h=1\,\mathrm{kWh}$. $100\,000\times700\,\mathrm W=70\,\mathrm{MW}$; en 24 h: $\mathbf{1\,680\ MWh}$ (solo GPU, sin refrigeración).
\textbf{(b)} $25(0.3)=7.5$ Wh por tarea; por tarea correcta $7.5/0.75=\mathbf{10\ Wh}$ (los intentos fallidos también consumen).
\textbf{(c)} $0.4\times3=1.2$: el consumo total \textbf{sube 20\,\%} pese a la eficiencia (efecto rebote).
\end{solucion}

% ============================================================
\newpage
\section{Parte C — Problemas tipo examen}

\begin{solucion}{C1 · Diagnóstico de un servidor}
\textbf{(a)} $P(A)P(U)P(C\mid A,U)P(M\mid C)P(Q\mid C)$: $1+1+4+2+2=\mathbf{10}$ parámetros frente a $2^5-1=31$.

\textbf{(b)}
$A\indep U$: \textbf{sí} ($A\to C\leftarrow U$, colisionador sin evidencia).
$A\indep U\mid M$: \textbf{no} ($M$ es descendiente de $C$: abre el colisionador).
$M\indep Q$: \textbf{no} ($M\leftarrow C\to Q$, causa común no observada).
$M\indep Q\mid C$: \textbf{sí} (causa común observada).
$A\indep Q\mid C$: \textbf{sí} (cadena $A\to C\to Q$ bloqueada).

\textbf{(c)} Factores tras restringir: $P(A)$, $P(U)$, $P(C\mid A,U)$, $P(+m\mid C)=(0.9,\,0.1)$. $Q$ es hoja no observada: se descarta.

\emph{Eliminar $C$}: $f_1(A,U)=\sum_cP(c\mid A,U)P(+m\mid c)$:
\begin{center}\small
\begin{tabular}{cc|l}
$A$ & $U$ & $f_1(A,U)$\\\hline
$+a$ & $+u$ & $0.95(0.9)+0.05(0.1)=0.86$\\
$+a$ & $-u$ & $0.80(0.9)+0.20(0.1)=0.74$\\
$-a$ & $+u$ & $0.60(0.9)+0.40(0.1)=0.58$\\
$-a$ & $-u$ & $0.05(0.9)+0.95(0.1)=0.14$
\end{tabular}
\end{center}
\emph{Eliminar $U$}: $f_2(A)=\sum_uP(u)f_1(A,u)$: $f_2(+a)=0.1(0.86)+0.9(0.74)=0.752$; $f_2(-a)=0.1(0.58)+0.9(0.14)=0.184$.

\emph{Unir con $P(A)$ y normalizar}: $0.05(0.752)=0.0376$; $0.95(0.184)=0.1748$; $P(+m)=0.2124$.
\[P(+a\mid+m)=\frac{0.0376}{0.2124}=\mathbf{0.1770}.\]

\textbf{(d)} Con $+u$ no se suma $U$: pesos $P(A)P(+u)f_1(A,+u)$: $0.05(0.1)(0.86)=0.0043$ y $0.95(0.1)(0.58)=0.0551$.
$P(+a\mid+m,+u)=0.0043/0.0594=\mathbf{0.0724}$. La actualización defectuosa ya explica la CPU alta y la alerta: la creencia de ataque cae de
17.7\,\% a 7.2\,\% (\emph{explaining away}; recuerda que $A\indep U$ sin evidencia, pero no dada $M$).

\textbf{(e)} Ahora cada término incluye $P(+q\mid c)$: $g_1(A,U)=\sum_cP(c\mid A,U)P(+m\mid c)P(+q\mid c)$, con $P(+m\mid c)P(+q\mid c)=0.63$ ($+c$) y $0.005$ ($-c$):
$g_1=0.59875,\ 0.505,\ 0.38,\ 0.03625$ para $(+a,+u),(+a,-u),(-a,+u),(-a,-u)$.
$g_2(+a)=0.1(0.59875)+0.9(0.505)=0.514375$; $g_2(-a)=0.1(0.38)+0.9(0.03625)=0.070625$.
Pesos $0.05(0.514375)=0.02572$ y $0.95(0.070625)=0.06709$: $P(+a\mid+m,+q)=\mathbf{0.2771}$.
$M$ y $Q$ solo informan sobre $C$. Aun sabiendo con certeza $+c$, $P(+a\mid+c)=0.2900$: ese es el techo; la segunda señal acerca $C$ a la certeza
pero no distingue ataque de actualización.
\end{solucion}

\begin{solucion}{C2 · Enumeración frente a eliminación}
\textbf{(a)} Mundos $P(d)P(i)P(c\mid d,i)P(+e\mid i)P(+r\mid c)$:
\begin{center}\small
\begin{tabular}{ccc|l@{\qquad}ccc|l}
$I$&$D$&$C$& & $I$&$D$&$C$& \\\hline
$+i$&$+d$&$+c$& $0.4(0.3)(0.7)(0.8)(0.9)=0.06048$ & $-i$&$+d$&$+c$& $0.4(0.7)(0.1)(0.1)(0.9)=0.00252$\\
$+i$&$+d$&$-c$& $0.4(0.3)(0.3)(0.8)(0.2)=0.00576$ & $-i$&$+d$&$-c$& $0.4(0.7)(0.9)(0.1)(0.2)=0.00504$\\
$+i$&$-d$&$+c$& $0.6(0.3)(0.9)(0.8)(0.9)=0.11664$ & $-i$&$-d$&$+c$& $0.6(0.7)(0.4)(0.1)(0.9)=0.01512$\\
$+i$&$-d$&$-c$& $0.6(0.3)(0.1)(0.8)(0.2)=0.00288$ & $-i$&$-d$&$-c$& $0.6(0.7)(0.6)(0.1)(0.2)=0.00504$
\end{tabular}
\end{center}
Sumas: $0.18576$ y $0.02772$; $P(+r,+e)=0.21348$; $P(+i\mid+r,+e)=\mathbf{0.8702}$.

\textbf{(b)} Los factores $f_1(C,I)$ y $f_2(I)=(0.774,\,0.396)$ de B8 no cambian; solo se añade $P(+e\mid I)=(0.8,\,0.1)$ al final:
$0.3(0.8)(0.774)=0.18576$ y $0.7(0.1)(0.396)=0.02772$ $\Rightarrow$ $\mathbf{0.8702}$.

\textbf{(c)} Enumeración: $2^2=4$ mundos por valor de $I$, 8 productos de 5 factores (32 multiplicaciones). Eliminación: el mayor producto es
$P(D)P(C\mid D,I)$ con alcance $\{C,D,I\}$: 8 entradas. En redes grandes la diferencia es exponencial.

\textbf{(d)} $P(+i\mid+c)=\frac{0.246}{0.442}=\mathbf{0.5566}$; $P(+i\mid+c,+d)=\frac{0.084}{0.112}=\mathbf{0.75}$; $P(+i\mid+c,-d)=\frac{0.162}{0.330}=\mathbf{0.4909}$.
Un curso difícil hace \emph{menos} probable una nota alta; si aun así la nota fue alta, hace falta una explicación más fuerte: inteligencia.
No contradice el \emph{explaining away}: ambos son razonamiento intercausal a través del colisionador $C$ observado. Cuando las dos causas
\emph{favorecen} el efecto, confirmar una baja la otra; cuando una lo \emph{desfavorece} (como $D$), confirmarla sube la otra.
\end{solucion}

\begin{solucion}{C3 · Comparar métodos de muestreo}
\textbf{(a)} Con $+m$: muestras 2 $(-a,+u,+c)$, 4 $(+a,-u,+c)$, 6 $(-a,-u,-c)$, 12 $(-a,+u,+c)$. Se aceptan \textbf{4} y una tiene $+a$: $\hat P=\mathbf{0.25}$
(exacto 0.1770). Esperadas: $12\,P(+m)=12(0.2124)\approx\mathbf{2.5}$; aquí tocaron más.

\textbf{(b)} Peso $=P(+m\mid c)$: $0.9$ si $+c$, $0.1$ si $-c$.
Pesos: $0.1,0.9,\mathbf{0.9},0.1,0.9,0.1,0.1,0.1$ (en negrita la única con $+a$). Suma $3.2$; $\hat P(+a\mid+m)=0.9/3.2=\mathbf{0.2813}$.

\textbf{(c)} $P(A\mid+u,+c)\propto P(A)P(+c\mid A,+u)$: $+a$: $0.05(0.95)=0.0475$; $-a$: $0.95(0.6)=0.57$; $P(+a\mid\cdot)=\mathbf{0.0769}$.
Con $A=-a$: $P(C\mid-a,+u,+m)\propto P(C\mid-a,+u)P(+m\mid C)$: $+c$: $0.6(0.9)=0.54$; $-c$: $0.4(0.1)=0.04$; $P(+c\mid\cdot)=\mathbf{0.9310}$.

\textbf{(d)} Rechazo descartó 8 de 12 muestras (desperdicio $\approx1-P(+m)$, es decir 79\,\%). Verosimilitud con pesos usa todas, pero con pesos desiguales.
Gibbs no desperdicia muestras, pero están correlacionadas. Con evidencia $+m,+q$ ($P=0.0928$, y menor aún con más evidencia) rechazo es el peor;
preferiría verosimilitud con pesos (la evidencia es hoja: pesos no tan extremos) o Gibbs.
\end{solucion}

\begin{solucion}{C4 · Clasificador bayesiano ingenuo}
\textbf{(a)} $J\to\text{Cielo}$, $J\to\text{Viento}$: los atributos son independientes dada la clase.
$\hat j=\argmax_jP(j)P(\text{cielo}\mid j)P(\text{viento}\mid j)$.

\textbf{(b)} $P(J=\text{sí})=6/10$. Clase sí (6): cielo sol 2, nub 3, lluvia 1; viento sí 1, no 5. Clase no (4): sol 2, nub 0, lluvia 2; viento sí 4, no 0.
\begin{center}\small
\begin{tabular}{c|ccc|cc}
 & sol & nub & lluvia & viento sí & viento no\\\hline
$J=$sí & 2/6 & 3/6 & 1/6 & 1/6 & 5/6\\
$J=$no & 2/4 & 0 & 2/4 & 4/4 & 0
\end{tabular}
\end{center}

\textbf{(c)} (lluvia, sí): sí $0.6\cdot\frac16\cdot\frac16=0.0167$; no $0.4\cdot\frac12\cdot1=0.2$ $\Rightarrow P(\text{sí})=0.0769$: \textbf{no}.
(nub, sí): sí $0.6\cdot\frac12\cdot\frac16=0.05$; no $0.4\cdot0\cdot1=0$ $\Rightarrow P(\text{sí})=1$: \textbf{sí}.
(sol, no): sí $0.6\cdot\frac13\cdot\frac56=0.1667$; no $0.4\cdot\frac12\cdot0=0$ $\Rightarrow1$: \textbf{sí}.
Problema: un solo conteo cero anula todo el producto y da certeza absoluta a partir de 4 ejemplos.

\textbf{(d)} Laplace (cielo $k=3$, viento $k=2$):
\begin{center}\small
\begin{tabular}{c|ccc|cc}
 & sol & nub & lluvia & viento sí & viento no\\\hline
$J=$sí & 3/9 & 4/9 & 2/9 & 2/8 & 6/8\\
$J=$no & 3/7 & 1/7 & 3/7 & 5/6 & 1/6
\end{tabular}
\end{center}
(lluvia, sí): $0.6\cdot\frac29\cdot\frac28=0.0333$ vs $0.4\cdot\frac37\cdot\frac56=0.1429$ $\Rightarrow P(\text{sí})=\mathbf{0.1892}$ (no).
(nub, sí): $0.0667$ vs $0.0476$ $\Rightarrow\mathbf{0.5833}$ (sí).
(sol, no): $0.15$ vs $0.0286$ $\Rightarrow\mathbf{0.84}$ (sí).
Las decisiones no cambian; cambia la confianza: ya no hay certezas de 1.
\end{solucion}

\begin{solucion}{C5 · EM con dos monedas}
\textbf{(a)} $Z\in\{A,B\}$ oculta con $P(Z)=\frac12$; $Z\to H$ (caras en 10). $L(\theta)=\prod_j\big[\tfrac12\theta_A^{h_j}(1-\theta_A)^{10-h_j}+\tfrac12\theta_B^{h_j}(1-\theta_B)^{10-h_j}\big]$
(el coeficiente binomial es común y se omite).

\textbf{(b)} Ejemplo $h=5$: $0.6^5\,0.4^5=7.96\times10^{-4}$, $0.5^{10}=9.77\times10^{-4}$, $w=0.4491$.
\begin{center}\small
\begin{tabular}{c|ccccc}
$h_j$ & 5 & 9 & 8 & 4 & 7\\\hline
$w_j=P(A\mid h_j)$ & 0.4491 & 0.8050 & 0.7335 & 0.3522 & 0.6472\\
$1-w_j$ & 0.5509 & 0.1950 & 0.2665 & 0.6478 & 0.3528
\end{tabular}
\end{center}

\textbf{(c)} Caras esperadas de $A$: $\sum w_jh_j=21.30$; cruces: $\sum w_j(10-h_j)=8.57$. De $B$: $33-21.30=11.70$ caras y $17-8.57=8.43$ cruces.
$\theta_A=21.30/29.87=\mathbf{0.7130}$; $\theta_B=11.70/20.13=\mathbf{0.5813}$.
La log-verosimilitud pasa de $-33.09$ a $-31.86$. (Una segunda iteración da $\theta_A=0.7453$, $\theta_B=0.5693$.)

\textbf{(d)} Con la moneda conocida: MLE directo, $\theta_A=$ caras de $A$ / lanzamientos de $A$. Intercambiar las etiquetas da la misma
verosimilitud: la solución es simétrica (no identificable); la inicialización elige cuál es «$A$».
\end{solucion}

\begin{solucion}{C6 · Regresión polinomial y generalización}
\textbf{(a)} $X^\mathsf TX=\begin{pmatrix}5&0\\0&10\end{pmatrix}$, $X^\mathsf T\mathbf y=(10,\,-0.5)$; $\hat\theta=(\mathbf{2},\mathbf{-0.05})$.
Predicciones $2.1,2.05,2,1.95,1.9$; residuos $2.1,-1.25,-1.9,-0.85,1.9$; ECM$_{\text{tr}}=13.915/5=\mathbf{2.783}$.
Prueba: $\hat y(3)=1.85$, $\hat y(-3)=2.15$; ECM$_{\text{te}}=(7.25^2+6.65^2)/2=\mathbf{48.39}$.

\textbf{(b)} Con $\phi=(1,x,x^2)$: $\sum x=\sum x^3=0$, $\sum x^2=10$, $\sum x^4=34$, $\sum x^2y=33.9$:
\[X^\mathsf TX=\begin{pmatrix}5&0&10\\0&10&0\\10&0&34\end{pmatrix},\quad X^\mathsf T\mathbf y=\begin{pmatrix}10\\-0.5\\33.9\end{pmatrix}.\]
La fila central da $\theta_1=-0.05$. Las otras: $5\theta_0+10\theta_2=10$, $10\theta_0+34\theta_2=33.9$; $\det=70$:
$\theta_0=\frac{340-339}{70}=\mathbf{0.0143}$, $\theta_2=\frac{169.5-100}{70}=\mathbf{0.9929}$.
Residuos de entrenamiento $0.114,-0.257,0.086,0.143,-0.086$: ECM$_{\text{tr}}=\mathbf{0.0229}$.
Prueba: $\hat y(3)=8.80$, $\hat y(-3)=9.10$; errores $\pm0.3$: ECM$_{\text{te}}=\mathbf{0.09}$.

\textbf{(c)} En los 5 puntos el polinomio reproduce $y$ exactamente (p.\,ej.\ $x=1$: $0.1+0.2333+0.8083-0.0833+0.0417=1.1$): ECM$_{\text{tr}}=0$.
$\hat y(3)=0.1+0.7+7.275-2.25+3.375=9.2$; $\hat y(-3)=0.1-0.7+7.275+2.25+3.375=12.3$; ECM$_{\text{te}}=(0.01+12.25)/2=\mathbf{6.13}$.

\textbf{(d)} Entrenamiento: grado 4 (0) $<$ grado 2 (0.023) $<$ grado 1 (2.78). Prueba: grado 2 (0.09) $<$ grado 4 (6.13) $<$ grado 1 (48.4).
Grado 1: \textbf{subajuste}; grado 4: \textbf{sobreajuste} (sigue el ruido con 5 parámetros y 5 datos); grado 2 coincide con la estructura real.
\end{solucion}

\begin{solucion}{C7 · $k$-means y escala}
\textbf{(a)} Crudo, $\mu_1=(10,1)$, $\mu_2=(12,9)$. Iteración 1: $d^2(C_3)=(401,\,373)$, $d^2(C_4)=(533,\,401)$ $\Rightarrow$ $\{C_1\}$, $\{C_2,C_3,C_4\}$;
$\mu_2=(24.67,6.33)$, $J=271.3$. Iteración 2: $C_2$: $d^2=(68,\,167.6)$ pasa al grupo 1; $C_3$: $(401,\,47.2)$; $C_4$: $(533,\,56.6)$ $\Rightarrow$
$\{C_1,C_2\}$, $\{C_3,C_4\}$; $\mu_1=(11,5)$, $\mu_2=(31,5)$, $J=17+17+10+10=\mathbf{54}$. Iteración 3 sin cambios. \textbf{Agrupa por ingreso.}

\textbf{(b)} Escala: $C_1(1,1)$, $C_2(1.2,9)$, $C_3(3,2)$, $C_4(3.2,8)$. $C_3$: $d^2=(5,\,52.24)$; $C_4$: $(53.84,\,5)$ $\Rightarrow$ $\{C_1,C_3\}$, $\{C_2,C_4\}$;
$\mu_1=(2,1.5)$, $\mu_2=(2.2,8.5)$, $J=4(1.25)=\mathbf{5}$; estable. \textbf{Agrupa por visitas.} (La estandarización $z$ da la misma partición.)

\textbf{(c)} Crudo: $a(C_1)=d(C_1,C_2)=8.246$, $b=\frac{20.025+23.087}2=21.556$, $s=\mathbf{0.617}$.
Escalado: $a=d(C_1,C_3)=2.236$, $b=\frac{8.002+7.338}2=7.670$, $s=\mathbf{0.708}$.

\textbf{(d)} Ninguno es «el correcto» en abstracto: los dos son óptimos locales de criterios distintos. Lo decide el objetivo (¿segmentar por
poder adquisitivo o por frecuencia?) y, por tanto, la elección de escala; la silueta solo compara dentro de una misma escala.
\end{solucion}

\begin{solucion}{C8 · Robot de reciclaje}
\textbf{(a)}
\begin{center}\small
\begin{tabular}{cc|ll}
\toprule
$s$ & $a$ & $s'$ (prob.) & $R$\\\midrule
$H$ & buscar & $H$ (0.7), $L$ (0.3) & $+2$, $+2$\\
$H$ & esperar & $H$ (1) & $+1$\\
$L$ & buscar & $L$ (0.6), $H$ (0.4) & $+2$, $-3$\\
$L$ & esperar & $L$ (1) & $+1$\\
$L$ & recargar & $H$ (1) & $0$\\\bottomrule
\end{tabular}
\end{center}

\textbf{(b)} \emph{Iteración de políticas.}
\begin{itemize}
\item $\pi_0=(\text{esp},\text{esp})$: $V=1+0.9V\Rightarrow V(H)=V(L)=10$. Mejora: $H$: buscar $0.7(2+9)+0.3(2+9)=11>10$; $L$: buscar $0.6(11)+0.4(-3+9)=9$,
esperar 10, recargar 9 $\Rightarrow$ $\pi_1=(\text{buscar},\text{esp})$.
\item $\pi_1$: $V(L)=10$; $V(H)=2+0.63V(H)+0.27(10)\Rightarrow V(H)=4.7/0.37=12.7027$. Mejora: $H$: buscar 12.7027 vs esperar 12.4324;
$L$: buscar $6.6+0.4(-3+11.4324)=9.9730$, esperar 10, recargar $0.9(12.7027)=11.4324$ $\Rightarrow$ $\pi_2=(\text{buscar},\text{recargar})$.
\item $\pi_2$: $V(L)=0.9V(H)$, $V(H)=2+0.63V(H)+0.27(0.9V(H))\Rightarrow0.127V(H)=2$: $V(H)=\mathbf{15.748}$, $V(L)=\mathbf{14.173}$. Mejora: $H$: buscar 15.748
vs esperar 15.173; $L$: buscar 13.323, esperar 13.756, recargar 14.173. Sin cambios: \textbf{$\pi^*=(\text{buscar},\text{recargar})$}.
\end{itemize}

\textbf{(c)} \emph{Iteración de valores.}
\begin{center}\small
\begin{tabular}{c|cc|ccc|c}
\toprule
$k$ & $H$:buscar & $H$:esperar & $L$:buscar & $L$:esperar & $L$:recargar & $(V_k(H),V_k(L))$\\\midrule
1 & 2 & 1 & 0 & 1 & 0 & $(2,1)$\\
2 & 3.53 & 2.8 & 1.26 & 1.9 & 1.8 & $(3.53,1.9)$\\
3 & 4.7369 & 4.177 & 2.2968 & 2.71 & 3.177 & $(4.7369,3.177)$\\
4 & 5.842 & 5.2632 & 3.4209 & 3.8593 & 4.2632 & $(5.842,4.2632)$\\\bottomrule
\end{tabular}
\end{center}
Ej.\ $k=2$, $H$:buscar $=0.7(2+0.9\cdot2)+0.3(2+0.9\cdot1)=3.53$. Con $V_4$ la voraz es (buscar, recargar), la óptima; con $V_1,V_2$ en $L$ ganaba esperar
(horizonte corto). Los valores aún están lejos de $(15.75,14.17)$: convergen geométricamente con razón $0.9$.

\textbf{(d)} Solo cambia $Q(L,\text{buscar})$, que empeora; la política óptima no usa esa acción, así que $\pi^*$ y $V^*$ \textbf{no cambian}.
\end{solucion}

\begin{solucion}{C9 · Q-learning en el corredor}
\textbf{(a)} $Q\leftarrow Q+0.5[y-Q]$; terminales $L,R$ con $y=r$.
\begin{center}\small
\begin{tabular}{c|c|c|c}
Ep. & Transición & $y$ & $Q$ anterior $\to$ nuevo\\\hline
1 & $(2,\text{der},-1,3)$ & $-1+0.9(0)=-1$ & $Q(2,\text{der})$: $0\to-0.5$\\
1 & $(3,\text{der},+10,R)$ & $10$ & $Q(3,\text{der})$: $0\to5$\\
2 & $(1,\text{der},-1,2)$ & $-1+0.9\max(0,-0.5)=-1$ & $Q(1,\text{der})$: $0\to-0.5$\\
2 & $(2,\text{der},-1,3)$ & $-1+0.9(5)=3.5$ & $Q(2,\text{der})$: $-0.5\to1.5$\\
2 & $(3,\text{izq},-1,2)$ & $-1+0.9(1.5)=0.35$ & $Q(3,\text{izq})$: $0\to0.175$\\
2 & $(2,\text{der},-1,3)$ & $-1+0.9(5)=3.5$ & $Q(2,\text{der})$: $1.5\to2.5$\\
2 & $(3,\text{der},+10,R)$ & $10$ & $Q(3,\text{der})$: $5\to7.5$\\
3 & $(2,\text{izq},-1,1)$ & $-1+0.9\max(0,-0.5)=-1$ & $Q(2,\text{izq})$: $0\to-0.5$\\
3 & $(1,\text{izq},+2,L)$ & $2$ & $Q(1,\text{izq})$: $0\to1$\\
4 & $(1,\text{der},-1,2)$ & $-1+0.9(2.5)=1.25$ & $Q(1,\text{der})$: $-0.5\to0.375$\\
4 & $(2,\text{der},-1,3)$ & $-1+0.9(7.5)=5.75$ & $Q(2,\text{der})$: $2.5\to4.125$\\
4 & $(3,\text{der},+10,R)$ & $10$ & $Q(3,\text{der})$: $7.5\to8.75$
\end{tabular}
\end{center}
Tabla final: $Q(1,\cdot)=(1,\ 0.375)$, $Q(2,\cdot)=(-0.5,\ 4.125)$, $Q(3,\cdot)=(0.175,\ 8.75)$ (izq, der).

\textbf{(b)} Voraz: 1 $\to$ \textbf{izq}, 2 $\to$ der, 3 $\to$ der.

\textbf{(c)} $Q^*$ (de B24): $(2.504,\,4.016)$, $(3.495,\,6.134)$, $(5.616,\,8.904)$. El estado 1 está equivocado: $Q(1,\text{der})$ se actualizó cuando
$Q(2,\cdot)$ aún era bajo, y $Q(1,\text{izq})$ recibió de golpe medio $+2$. Hace falta más exploración (y más episodios) para que el valor de la
derecha se propague hasta el estado 1.

\textbf{(d)} El objetivo de Q-learning usa $\max_{a'}Q(s',a')$, no la acción que el agente tomará después: la exploración solo decide qué
entrada se actualiza, sin contaminar los objetivos de las demás. Con SARSA, la transición $(2,\text{der},-1,3)$ seguida de «izq» habría usado
$Q(3,\text{izq})=0$: $y=-1$ y $Q(2,\text{der})=-0.75$ en vez de $1.5$. Si hubiera pedido «der» y resbalado a 2, la tupla sería
$(3,\text{der},-1,2)$ y se actualizaría $Q(3,\text{der})$: el ruido del entorno queda promediado dentro de ese valor.
\end{solucion}

\begin{solucion}{C10 · Monte Carlo frente a TD}
\textbf{(a)} $V(A)$: un solo retorno, $0+0=0$ $\Rightarrow$ $\mathbf0$. $V(B)$: retornos $0,1,1,1,1,1,1,0$ $\Rightarrow$ $6/8=\mathbf{0.75}$.
\textbf{(b)} En el punto fijo batch, cada valor es el promedio de sus objetivos. $B$ siempre termina: $V(B)=0.75$. La única transición de $A$
va a $B$ con 0: $V(A)=0+V(B)=\mathbf{0.75}$.
\textbf{(c)} TD es la respuesta del \emph{modelo de máxima verosimilitud} (A siempre pasa a B con 0, y B paga 1 con prob.\ 3/4) suponiendo Markov;
MC minimiza el error cuadrático sobre los retornos observados (y da 0 por un único episodio). Si el proceso es Markov, TD predice mejor el futuro; MC
no supone Markov.
\textbf{(d)} Ep.\,1: $V(A)\leftarrow0+0.5(0+0-0)=0$; $V(B)\leftarrow0$. Luego seis veces $B\to1$: $0.5,\,0.75,\,0.875,\,0.9375,\,0.96875,\,0.984375$; el último
$B\to0$: $0.4922$. Resultado $V(A)=\mathbf0$, $V(B)=\mathbf{0.4922}$: en línea con $\alpha$ fija, pesan más las experiencias recientes.
\end{solucion}

\begin{solucion}{C11 · Retropropagación con softmax}
\textbf{(a)} $\mathbf z_1=W_1\mathbf x+\mathbf b_1=(2-1+0,\ 1+1-3)=(1,-1)$; $\mathbf h=(1,0)$; $\mathbf z_2=W_2\mathbf h=(1,\,0.5)$;
$\mathbf p=\mathrm{softmax}(1,0.5)=(0.6225,\,0.3775)$; $\ell=-\ln0.3775=\mathbf{0.9741}$.

\textbf{(b)} $\boldsymbol\delta_2=\mathbf p-\mathbf t=(0.6225,-0.6225)$;
$\frac{\partial\ell}{\partial W_2}=\boldsymbol\delta_2\mathbf h^\mathsf T=\begin{pmatrix}0.6225&0\\-0.6225&0\end{pmatrix}$, $\frac{\partial\ell}{\partial\mathbf b_2}=\boldsymbol\delta_2$.
$W_2^\mathsf T\boldsymbol\delta_2=(1(0.6225)+0.5(-0.6225),\ -1(0.6225)+1(-0.6225))=(0.3113,\,-1.2449)$.
$\boldsymbol\delta_1=(0.3113,\,-1.2449)\odot(1,0)=(0.3113,\,0)$;
$\frac{\partial\ell}{\partial W_1}=\boldsymbol\delta_1\mathbf x^\mathsf T=\begin{pmatrix}0.3113&-0.3113\\0&0\end{pmatrix}$, $\frac{\partial\ell}{\partial\mathbf b_1}=(0.3113,0)$.

\textbf{(c)} $\eta=1$: $W_1=\begin{pmatrix}1.6888&1.3112\\1&-1\end{pmatrix}$, $\mathbf b_1=(-0.3113,-3)$, $W_2=\begin{pmatrix}0.3775&-1\\1.1225&1\end{pmatrix}$, $\mathbf b_2=(-0.6225,0.6225)$.
Nuevo forward: $\mathbf z_1=(0.0663,-1)$, $\mathbf h=(0.0663,0)$, $\mathbf z_2=(-0.5974,\,0.6969)$, $\mathbf p=(0.2151,\,0.7849)$, $\ell=\mathbf{0.2422}$.

\textbf{(d)} La segunda unidad oculta está apagada ($z=-1$): no cambian la segunda fila de $W_1$, $b_{1,2}$, ni la segunda columna de $W_2$
(multiplica $h_2=0$). Neurona muerta: si está apagada para \emph{todas} las entradas, nunca recibe gradiente.
\end{solucion}

\begin{solucion}{C12 · Un transformer diminuto}
\textbf{(a)}
\begin{center}\small
\begin{tabular}{l|l|c|c}
Prefijo & Cálculo & $p(\text{correcto})$ & $-\ln p$\\\hline
$\langle s\rangle$ & $e^2/(e^2+3+e^{-1})=7.389/10.757$ & 0.6869 & 0.3756\\
… el & $e^2/(e^2+e+3)=7.389/13.107$ & 0.5637 & 0.5732\\
… gato & $e^2/(e+1+e^2+1+1)$ & 0.5637 & 0.5732\\
… come & $e^3/(1+1+e+1+e^3)=20.09/25.80$ & 0.7784 & 0.2505
\end{tabular}
\end{center}
Pérdida media $=1.7724/4=\mathbf{0.4431}$; perplejidad $e^{0.4431}=\mathbf{1.558}$ (como elegir entre $\approx1.56$ opciones por paso).

\textbf{(b)} $\tau=0.5$ (logits $\times2$): $(0.1137,0.0154,\mathbf{0.8401},0.0154,0.0154)$. $\tau=2$ (logits $/2$): $(0.2238,0.1357,\mathbf{0.3690},0.1357,0.1357)$.

\textbf{(c)} $S=QK^\mathsf T/\sqrt2=\begin{pmatrix}0.7071&0.7071&1.4142\\0&1.4142&1.4142\\1.4142&0&1.4142\end{pmatrix}$. Con máscara:
$A=\begin{pmatrix}1&0&0\\0.1956&0.8044&0\\0.4458&0.1084&0.4458\end{pmatrix}$,
$O=AV=\begin{pmatrix}2&0\\0.3911&1.6089\\1.3374&0.6626\end{pmatrix}$.
(Fila 3: $0.4458(2,0)+0.1084(0,2)+0.4458(1,1)$.)

\textbf{(d)} Embeddings: $50\,000\times512=\mathbf{25\,600\,000}$. $W_Q,W_K,W_V$: $3\times512\times64=98\,304$; $W_O$: $64\times512=32\,768$; atención: $\mathbf{131\,072}$.
El vocabulario domina en modelos pequeños.
\end{solucion}

\begin{solucion}{C13 · Auditoría de un sistema de crédito}
\textbf{(a)} Exactitud global $(48+32+18+68)/200=\mathbf{0.83}$.
\begin{center}\small
\begin{tabular}{c|cccccc}
Grupo & Tasa base & Aprobación & TPR & FPR & PPV & Exactitud\\\hline
A & 0.60 & 0.56 & 0.80 & 0.20 & 0.857 & 0.80\\
B & 0.30 & 0.20 & 0.60 & 0.029 & 0.90 & 0.86
\end{tabular}
\end{center}
\textbf{(b)} Se violan paridad ($0.56$ vs $0.20$), igualdad de oportunidades ($0.80$ vs $0.60$) y \emph{equalized odds}. En B, el 40\,\% de quienes pagarían
es rechazado (costo para esas personas); en A, el 20\,\% de quienes no pagarían es aprobado (costo para el prestamista y riesgo de
sobreendeudamiento). La mejor exactitud de B no compensa esa asimetría.
\textbf{(c)} Las tasas base difieren (0.6 vs 0.3): con PPV (calibración) igual y tasas base distintas, no pueden igualarse también FPR y FNR salvo con un
clasificador perfecto. Etiquetas selectivas: solo se observa si paga quien recibió crédito; los rechazados nunca generan la etiqueta, y el
sesgo se perpetúa en el siguiente entrenamiento.
\textbf{(d)} Alto riesgo (evaluación de solvencia de personas físicas, anexo III): gestión de riesgos, gobernanza de datos, documentación técnica,
registros, transparencia, supervisión humana, exactitud y robustez, y evaluación de conformidad antes de ponerlo en el mercado.
\textbf{(e)} $40(0.3)=12$ Wh por solicitud; por solicitud bien resuelta $12/0.8=\mathbf{15\ Wh}$; $10\,000$ solicitudes diarias: $\mathbf{120\ kWh/día}$.
\end{solucion}

% ============================================================
\newpage
\section{Parte D — Simulacros resueltos}

\subsection*{Simulacro D1}
\addcontentsline{toc}{subsection}{Simulacro D1}

\begin{solucion}{D1.1 · El auto no arranca}
\textbf{(a)} $P(B)P(F)P(S\mid B,F)P(L\mid B)$: $1+1+4+2=\mathbf8$ parámetros (conjunta: 15).

\textbf{(b)} $B\indep F$: \textbf{sí} ($B\to S\leftarrow F$ sin evidencia). $B\indep F\mid S$: \textbf{no} (colisionador observado).
$L\indep F$: \textbf{sí} ($L\leftarrow B\to S\leftarrow F$ bloqueado en $S$). $L\indep F\mid S$: \textbf{no} ($S$ abre y $B$, bifurcación no observada, deja pasar).

\textbf{(c)} Pesos $P(b)P(f)P(-s\mid b,f)$:
\begin{center}\small
\begin{tabular}{cc|l}
$B$ & $F$ & $P(b,f,-s)$\\\hline
$+b$ & $+f$ & $0.9(0.95)(0.01)=0.00855$\\
$+b$ & $-f$ & $0.9(0.05)(1)=0.045$\\
$-b$ & $+f$ & $0.1(0.95)(0.95)=0.09025$\\
$-b$ & $-f$ & $0.1(0.05)(1)=0.005$
\end{tabular}
\end{center}
$P(-s)=0.1488$. $P(-b\mid-s)=\frac{0.09025+0.005}{0.1488}=\mathbf{0.6401}$; \quad $P(-f\mid-s)=\frac{0.045+0.005}{0.1488}=\mathbf{0.3360}$.

\textbf{(d)} Multiplicar cada fila por $P(+l\mid b)$ ($0.95$ si $+b$, $0.1$ si $-b$): $0.0081225$, $0.04275$, $0.009025$, $0.0005$; suma $0.0603975$.
$P(-b\mid-s,+l)=\frac{0.009525}{0.0603975}=\mathbf{0.1577}$; \quad $P(-f\mid-s,+l)=\frac{0.04325}{0.0603975}=\mathbf{0.7161}$.
Las luces encendidas indican batería buena; como el auto no arranca, la culpa pasa a la gasolina (razonamiento intercausal vía $S$).
\end{solucion}

\begin{solucion}{D1.2 · Aprender y muestrear}
\textbf{(a)}
\begin{center}\small
\begin{tabular}{l|ccc}
Fila & Conteo & MLE & Laplace\\\hline
$P(+s\mid+b,+f)$ & 4/4 & 1 & 5/6 = 0.833\\
$P(+s\mid+b,-f)$ & 0/1 & 0 & 1/3 = 0.333\\
$P(+s\mid-b,+f)$ & 1/3 & 0.333 & 2/5 = 0.4\\
$P(+s\mid-b,-f)$ & 0/0 & \textbf{indefinida} & 1/2\\
$P(+l\mid+b)$ & 4/5 & 0.8 & 5/7 = 0.714\\
$P(+l\mid-b)$ & 0/3 & 0 & 1/5 = 0.2
\end{tabular}
\end{center}

\textbf{(b)} Peso $=P(-s\mid b,f)$: $(1,1)\to0.01$, $(0,1)\to0.95$, $(1,0)\to1$, $(1,1)\to0.01$, $(0,1)\to0.95$, $(1,1)\to0.01$.
Suma $2.93$; con $+b$: $0.01+1+0.01+0.01=1.03$. $\hat P(+b\mid-s)=\mathbf{0.3515}$ frente a $1-0.6401=\mathbf{0.3599}$ exacto.

\textbf{(c)} Solo se acepta una fracción $\approx P(-s)$, es decir 14.9\,\% de las muestras: para 100 aceptadas se necesitan $\approx100/0.1488\approx\mathbf{672}$.
Además $L$ y $F$ se generarían sin usar la evidencia; con evidencia más rara el costo crece en proporción inversa a $P(\mathbf e)$.
\end{solucion}

\begin{solucion}{D1.3 · Inventario}
\textbf{(a)} $s'=s+a-\min(s+a,d)$, $R=4\min(s+a,d)-a-0.5s'$:
\begin{center}\small
\begin{tabular}{cc|ll}
$s$ & $a$ & Demanda 1 (0.5) & Demanda 0 (0.5)\\\hline
0 & 0 & $s'=0$, $R=0$ & $s'=0$, $R=0$\\
0 & 1 & $s'=0$, $R=3$ & $s'=1$, $R=-1.5$\\
1 & 0 & $s'=0$, $R=4$ & $s'=1$, $R=-0.5$\\
1 & 1 & $s'=1$, $R=2.5$ & $s'=2$, $R=-2$\\
2 & 0 & $s'=1$, $R=3.5$ & $s'=2$, $R=-1$
\end{tabular}
\end{center}

\textbf{(b)}
\begin{center}\small
\begin{tabular}{c|cc|cc|c|c}
\toprule
$k$ & $Q(0,0)$ & $Q(0,1)$ & $Q(1,0)$ & $Q(1,1)$ & $Q(2,0)$ & $V_k=(V_0,V_1,V_2)$\\\midrule
1 & 0 & 0.75 & 1.75 & 0.25 & 1.25 & $(0.75,1.75,1.25)$\\
2 & 0.6 & 1.75 & 2.75 & 1.45 & 2.45 & $(1.75,2.75,2.45)$\\
3 & 1.4 & 2.55 & 3.55 & 2.33 & 3.33 & $(2.55,3.55,3.33)$\\\bottomrule
\end{tabular}
\end{center}
Ej.\ $k=2$: $Q(0,1)=0.5(3+0.8\cdot0.75)+0.5(-1.5+0.8\cdot1.75)=1.8-0.05=1.75$. Voraz de $V_3$: pedir en 0, no pedir en 1.

\textbf{(c)} $\pi$ (pedir si $s<2$): $V_0=0.75+0.4V_0+0.4V_1$, $V_1=0.25+0.4V_1+0.4V_2$, $V_2=1.25+0.4V_1+0.4V_2$ $\Rightarrow$ $V^\pi=(3.417,\,3.25,\,4.25)$.
$\pi'$ (pedir solo en 0): $V_0=0.75+0.4V_0+0.4V_1$, $V_1=1.75+0.4V_0+0.4V_1$, $V_2=1.25+0.4V_1+0.4V_2$ $\Rightarrow$ $V_1-V_0=1$, $V^{\pi'}=(\mathbf{5.75},\mathbf{6.75},\mathbf{6.583})$.
$\pi'$ domina en todos los estados. Mejora con $V^{\pi'}$: $Q(0,0)=4.6<Q(0,1)=5.75$; $Q(1,0)=6.75>Q(1,1)=5.583$: no cambia $\Rightarrow$ $\pi'$ \textbf{es óptima}.

\textbf{(d)} Con $\gamma=0$ solo cuenta la recompensa inmediata ($V_1$ de (b)): pedir en 0 ($0.75>0$) y no en 1 ($1.75>0.25$). Aquí coincide con la óptima, pero en
general un agente miope ignora el costo futuro del inventario.
\end{solucion}

\begin{solucion}{D1.4 · Q-learning tabular y aproximado}
\textbf{(a)}
\begin{center}\small
\begin{tabular}{c|c|c}
Transición & $y=r+0.8\max_{a'}Q(s',a')$ & Nuevo $Q=Q+0.25(y-Q)$\\\hline
$(s_1,b,+1,s_2)$ & $1+0.8(4)=4.2$ & $Q(s_1,b)=2+0.25(2.2)=\mathbf{2.55}$\\
$(s_2,a,-2,s_3)$ & $-2+0.8(3)=0.4$ & $Q(s_2,a)=4+0.25(-3.6)=\mathbf{3.1}$\\
$(s_3,b,+5,T)$ & $5$ (terminal) & $Q(s_3,b)=3+0.25(2)=\mathbf{3.5}$\\
$(s_1,a,0,s_3)$ & $0+0.8(3.5)=2.8$ & $Q(s_1,a)=1+0.25(1.8)=\mathbf{1.45}$
\end{tabular}
\end{center}
\textbf{(b)} En $s_1$ la voraz es $b$ antes ($2>1$) y después ($2.55>1.45$): $\mu(b)=0.95$, $\mu(a)=0.05$.
\textbf{(c)} $\hat Q(s,a)=0.5+3-2=1.5$; en $s'$: $0.5+1+0=1.5$ y $0.5+2-2=0.5$; $y=-1+0.8(1.5)=0.2$; $\delta=-1.3$;
$\mathbf w'=(0.5,1,-2)+0.1(-1.3)(1,3,1)=(\mathbf{0.37},\mathbf{0.61},\mathbf{-2.13})$.
\end{solucion}

\subsection*{Simulacro D2}
\addcontentsline{toc}{subsection}{Simulacro D2}

\begin{solucion}{D2.1 · Lluvia, aspersor y tráfico}
\textbf{(a)}
$S\indep T$: \textbf{no} ($S\leftarrow N\to R\to T$ activo).
$S\indep T\mid R$: \textbf{sí} (ese camino se bloquea en $R$; $S\to W\leftarrow R\to T$ se bloquea en $W$ no observado).
$S\indep T\mid W$: \textbf{no} ($S\to W\leftarrow R\to T$: $W$ abre, $R$ no observado).
$S\indep T\mid R,W$: \textbf{sí} ($R$ observado bloquea ambos caminos: es no colisionador en los dos).
$N\indep W\mid R,S$: \textbf{sí} (padres de $W$ observados). $T\indep W\mid R$: \textbf{sí}.

\textbf{(b)} Factores: $P(N)$, $P(R\mid N)$, $P(S\mid N)$, $P(+w\mid R,S)$; $T$ es hoja no observada y se descarta.
Eliminar $N$: $f_1(R,S)=\sum_nP(n)P(R\mid n)P(S\mid n)$:
\begin{center}\small
\begin{tabular}{cc|l}
$R$ & $S$ & $f_1$\\\hline
$+r$ & $+s$ & $0.4(0.7)(0.2)+0.6(0.1)(0.6)=0.092$\\
$+r$ & $-s$ & $0.4(0.7)(0.8)+0.6(0.1)(0.4)=0.248$\\
$-r$ & $+s$ & $0.4(0.3)(0.2)+0.6(0.9)(0.6)=0.348$\\
$-r$ & $-s$ & $0.4(0.3)(0.8)+0.6(0.9)(0.4)=0.312$
\end{tabular}
\end{center}
Eliminar $S$: $f_2(+r)=0.092(0.95)+0.248(0.85)=0.2982$; $f_2(-r)=0.348(0.8)+0.312(0.05)=0.294$.
$P(+r\mid+w)=0.2982/0.5922=\mathbf{0.5035}$.

\textbf{(c)} Sin sumar $S$: $+r$: $0.092(0.95)=0.0874$; $-r$: $0.348(0.8)=0.2784$; $P(+r\mid+w,+s)=\mathbf{0.2389}$. Baja porque el aspersor explica el pasto
mojado y, además, $+s$ sugiere cielo despejado ($N$), lo que hace menos probable la lluvia.
\end{solucion}

\begin{solucion}{D2.2 · Regresión, logística y métricas}
\textbf{(a)} $X^\mathsf TX=\begin{pmatrix}4&12\\12&46\end{pmatrix}$, $X^\mathsf T\mathbf y=(20,77)$, $\det=40$;
$\hat\theta=\frac1{40}(46\cdot20-12\cdot77,\ -12\cdot20+4\cdot77)=(\mathbf{-0.1},\mathbf{1.7})$.
$\hat y=(1.6,3.3,6.7,8.4)$, residuos $(0.4,-0.3,-0.7,0.6)$, SSE $=\mathbf{1.1}$; $\hat y(6)=\mathbf{10.1}$.

\textbf{(b)} Con $\theta=0$: $\pi_i=0.5$, $J=4\ln2=\mathbf{2.7726}$. $\boldsymbol\pi-\mathbf y=(-0.5,-0.5,0.5,0.5)$;
$\nabla J=(0,\ 1(-0.5)+3(-0.5)+(-1)(0.5)+(-2)(0.5))=(0,\mathbf{-3.5})$; $\theta_1=(0,1.75)$.
$\pi=(0.8520,0.9948,0.1480,0.0293)$; $J_1=\mathbf{0.3554}$. Los datos son separables: la pérdida puede acercarse a 0 haciendo $\theta_1\to\infty$; sin regularización
$\theta$ crece indefinidamente (no hay mínimo finito).

\textbf{(c)} Exactitud $170/200=\mathbf{0.85}$; precisión $30/40=\mathbf{0.75}$; recall $30/50=\mathbf{0.6}$; especificidad $140/150=\mathbf{0.933}$; $F_1=60/90=\mathbf{0.667}$.
\end{solucion}

\begin{solucion}{D2.3 · Agrupamiento}
\textbf{(a)} Inicio $A,B$: asignación $C$: $d^2=(9,13)\to1$; $D,E,F\to2$ $\Rightarrow$ $\{A,C\}$, $\{B,D,E,F\}$; $\mu_1=(1.5,0)$, $\mu_2=(5.5,5.75)$, $J=74.25$.
Siguiente: $B$: $(6.25,\,44.31)\to1$ $\Rightarrow$ $\{A,B,C\}$, $\{D,E,F\}$; $\mu_1=(1,0.667)$, $\mu_2=(7.333,7)$, $J=\mathbf{19.33}$; estable.
Inicio $A,D$: converge en una iteración a la misma partición y $J$. Una mala inicialización cuesta iteraciones (y en otros datos puede dejar un óptimo local peor).

\textbf{(b)} Manhattan: $AB=2$, $AC=3$, $BC=5$, $DE=1$, $DF=6$, $EF=5$, $CD=9$, $CE=10$, $BD=10$, $AD=12$, $AF=18$.
\begin{center}\small
\begin{tabular}{c|ll|ll}
Paso & Simple & altura & Completo & altura\\\hline
1 & $D+E$ & 1 & $D+E$ & 1\\
2 & $A+B$ & 2 & $A+B$ & 2\\
3 & $C+AB$ & $\min(3,5)=3$ & $C+AB$ & $\max(3,5)=5$\\
4 & $F+DE$ & $\min(6,5)=5$ & $F+DE$ & $\max(6,5)=6$\\
5 & $ABC+DEF$ & $9$ ($C$–$D$) & $ABC+DEF$ & $18$ ($A$–$F$)
\end{tabular}
\end{center}
\textbf{(c)} Ambos: \textbf{2 grupos}; el mayor salto de altura está antes de la última unión (5 $\to$ 9 y 6 $\to$ 18).
\end{solucion}

\begin{solucion}{D2.4 · Red neuronal, atención y equidad}
\textbf{(a)} Forward: $z_1=2-1+0=1$, $h=1$, $z_2=2-1=1$, $p=0.7311$, $\ell=-\ln(1-p)=\mathbf{1.3133}$.
Backward: $\delta_2=p-y=0.7311$; $\partial\ell/\partial w_2=0.7311$, $\partial\ell/\partial b_2=0.7311$; $\delta_1=\delta_2w_2\cdot1=1.4621$;
$\partial\ell/\partial\mathbf w_1=(2.9242,1.4621)$, $\partial\ell/\partial b_1=1.4621$.
$\eta=0.5$: $\mathbf w_1=(-0.4621,-1.7311)$, $b_1=-0.7311$, $w_2=1.6345$, $b_2=-1.3655$.
Nuevo: $z_1=-3.386$, $h=0$, $p=\sigma(-1.3655)=0.2033$, $\ell=\mathbf{0.2273}$.

\textbf{(b)} Puntajes $\mathbf q^\mathsf T\mathbf k_j/\sqrt2=(2,2,3)/\sqrt2=(1.4142,1.4142,2.1213)$; pesos $(0.2483,0.2483,0.5035)$;
salida $0.2483(1,0)+0.2483(0,1)+0.5035(2,2)=(\mathbf{1.2552},\mathbf{1.2552})$. Con máscara en la posición 2 solo cuentan $\mathbf k_1,\mathbf k_2$:
pesos $(0.5,0.5)$ y salida $(0.5,0.5)$.

\textbf{(c)}
\begin{center}\small
\begin{tabular}{c|ccccc}
Grupo & Aprobación & TPR & FPR & PPV & Exactitud\\\hline
H & 0.45 & 0.75 & 0.25 & 0.667 & 0.75\\
M & 0.25 & 0.50 & 0.083 & 0.80 & 0.75
\end{tabular}
\end{center}
Ambos grupos tienen la misma tasa base (0.4), así que las diferencias no se explican por ella: se violan paridad, igualdad de oportunidades
(la mitad de las mujeres que pagarían son rechazadas) y \emph{equalized odds}, con exactitud idéntica.
\end{solucion}


\end{document}
